\newpage
\chapter{线性空间}
\setcounter{section}{0}
\section{数域}
\subsection{数域}
\begin{formal}
\begin{definition}[数域与数环]\label{def:数域与数环}
    设$\mathbb{K} \subseteq \mathbb{C} $
    且$\mathbb{K} $中元的个数不少于$2$.
    如果$\mathbb{K} $中任意两个元素
    作加、减、乘、除（除数不为零）
    得到的结果仍然属于$\mathbb{K} $,
    那么$\mathbb{K} $称为一个数域.
   
    特别地,如果$\mathbb{K} $仅对加、
    减、乘封闭而对除法不封闭,则称$\mathbb{K} $为一个
    数环.
\end{definition}
\end{formal}
\begin{brown}
    \begin{example}
回顾
\[
\mathbb{N}
\subseteq \mathbb{Z} 
\subseteq \mathbb{Q} 
\subseteq \mathbb{R} 
\subseteq \mathbb{C} 
\]
这其中,$\mathbb{N}$不是数域或数环,
$\mathbb{Z} $是整数环,
$\mathbb{Q} $是有理数域,
$\mathbb{R} $是实数域,
$\mathbb{C} $是复数域.
    \end{example}\end{brown}
    \begin{green}\begin{lemma}[最小的数域]\label{lem:最小的数域}
容易证明$\mathbb{C} $包含无穷多个数域
,但任一数域$\mathbb{K} 
\subseteq \mathbb{C}$一定满足：
\setcounter{equation}{0}%重置计数器,方程下计数器为equation
\begin{enumerate}[label=\textup{(\arabic*)}]
    \item $0 \in \mathbb{K} ,
    1 \in \mathbb{K} ; $
    \item $\mathbb{Q} 
    \subseteq \mathbb{K}$
\end{enumerate}
不难得知$\mathbb{Q}$是最“小”的数域.
\end{lemma}\end{green}
\begin{red}\begin{remark}
    一个包含数域（环）的数集不一定是
一个数域或数环.
\end{remark}
\end{red}
\subsection{超越数与代数数}
对任意的$\xi \in \mathbb{C}$,要么是
代数数,要么是超越数.
\[
\xi \in \mathbb{C} 
\begin{cases*}
    \text{代数数；} \\
    \text{超越数.}
\end{cases*}
\]
\begin{formal}
\begin{definition}[代数数]\label{def:代数数}
    设$\alpha  \in \mathbb{C}$
    ,如果存在一个多项式
    \[
    f\left(x\right) = a_n x^n + 
    \cdots + a_1 x + a_0 \neq 0
    \]
    其中$a_i \neq 0\left(i = 0,\cdots
    ,n\right)$,使得$f\left(\alpha 
    \right) = 0$,则称$
    \alpha $是一个代数数.
\end{definition}
\end{formal}
\begin{formal}
\begin{definition}[超越数]\label{def:超越数}
    设$\alpha  \in \mathbb{C}$
    ,如果不存在一个多项式
    \[
    f\left(x\right) = a_n x^n + 
    \cdots + a_1 x + a_0 \neq 0
    \]
    其中$a_i \neq 0\left(i = 0,\cdots
    ,n\right)$,使得$f\left(\alpha 
    \right) = 0$,
    或者说$f\left(\alpha \right)= 0
    $成立当且仅当
    \[
    a_i \equiv 0\quad\left(i = 0,\cdots
    ,n\right)
    \]则称$
    \alpha $是一个超越数.
\end{definition}\end{formal}
\begin{red}
\begin{remark}
    从抽象代数层面说,
    域（\textup{field}）是数域
    和四则运算的推广,是环的一种.最著名的域就是复数域、实数域、有理数域,此外还有有理函数域、
    代数函数域、代数数域（与数域有一些不同）等.

    域上定义有加法和乘法,一般
    地,记作$\oplus$和$\wedge$,最小
    的伽罗瓦域是$\mathbb{F}_2=\left\{\overline{0},
    \overline{1}\right\}$
\end{remark}
\end{red}
\begin{red}
\begin{remark}
    因此,数域只是域的一种,在高等代数中,我们主要讨论数域,即复数域$\mathbb{C}$的子域.
\end{remark}
\end{red}
\newpage
\section{向量}
\begin{red}
\begin{remark}
    此处的向量还是比较接近中学所学到的向量,但
    实际上
    向量并不只是一组坐标,这将在线性空间中体现.
\end{remark}
\end{red}
\subsection{定义}
\begin{formal}
\begin{definition}[$n$维行（列）向量]\label{def:n维行（列）向量}
    设数域$\mathbb{K}$,
    $a_1,a_2,\cdots,a_n \in \mathbb{K}$
    ,则$1 \times n$矩阵
    \[
    \left(a_1,a_2,\cdots,a_n\right)
    \]
    为$\mathbb{K}$上的$n$维行向量
    ；
    $n \times 1$矩阵\[
    \left(a_1,a_2,\cdots,a_n\right)^{\prime}
    =
    \begin{pmatrix}
        a_1 \\
        a_2 \\
        \vdots\\
        a_n
    \end{pmatrix}
    \]
    为$\mathbb{K}$上的$n$维列向量.
\end{definition}
\end{formal}
\subsection{n维行（列）向量空间}
\begin{formal}
\begin{definition}[$n$维行（列）向量空间]\label{def:n维行（列）向量空间}
    \[
    \mathbb{K}_n = \left\{
        \left(a_1,a_2,\cdots,a_n\right)
        | a_i \in \mathbb{K},
        1\leqslant i\leqslant n
    \right\}
    \]
    称为$\mathbb{K}$上的$n$维行向量空间；
    \[
    \mathbb{K}^n = \left\{
        \left(a_1,a_2,\cdots,a_n\right)^{\prime}
        | a_i \in \mathbb{K},
        1\leqslant i\leqslant n
    \right\}
    \]
    称为$\mathbb{K}$上的$n$维列向量空间.
\end{definition}\end{formal}
\subsection{向量的运算规则}
\begin{formal}
\begin{definition}[零向量]\label{def:零向量}
    若一个向量所有元素均为零,称为一个零向量.
\end{definition}
\end{formal}
\begin{formal}
\begin{definition}[负向量]\label{def:负向量} 
    设$\bm{\alpha}=\left(a_1,a_2,\cdots,a_n\right)$,记$-\bm{\alpha}=
    \left(-a_1,-a_2,\cdots,-a_n\right)$,称$-\bm{\alpha}$为$\bm{\alpha}$的负向量.
\end{definition}
\end{formal}
\begin{formal}
\begin{proposition}[$n$维向量空间的若干性质]\label{pop:n维向量空间的若干性质}
设$\bm{\alpha} $、$\bm{\beta} $、$\bm{
    \gamma} \in
\mathbb{K}^n$,$k$、$l\in \mathbb{K}$
\begin{enumerate}[label=\textup{(\arabic*)}]
    \item 加法交换律：$\bm{\alpha } + \bm{\beta} =
    \bm{\beta} +\bm{\alpha }
    $
    \item 加法结合律：$\left(\bm{\alpha} +\bm{ \beta}
     \right) +
    \bm{\gamma} =\left(
        \bm{\alpha} + \bm{\gamma }
    \right) + \bm{\beta}  = \bm{\alpha
     } +
    \left(\bm{\beta} + \bm{\gamma}
     \right)
    $
    \item 零元存在：$\bm{\alpha }+ \bm{0} = \bm{\alpha }
    $
\item 负元存在：$\bm{\alpha }+\left(-\bm{\alpha }
    \right) = \bm{0}
    $
    \item 单位元存在：$1\cdot\bm{\alpha } = \bm{\alpha }
    $
    \item 数乘的数分配律$ \left(k + l\right)\bm{\alpha }
    =k\bm{\alpha }+l\bm{\alpha }
    $
    \item 数乘的向量分配律：$k\left(\bm{\alpha }+
    \bm{\beta }\right) = 
    k\bm{\alpha }+ k\bm{\beta }
    $
    \item 数乘结合律：$\left(kl\right)\bm{\alpha }
    =l\left(k\bm{\alpha }\right)
    =k\left(l\bm{\alpha }\right)$
\end{enumerate}
\end{proposition}
\end{formal}
\begin{red}
    \begin{remark}
        此处的向量空间本身就是特殊的线性空间\cref{def:线性空间},因此性质与线性空间的性质是一致的.
    \end{remark}
\end{red}
\newpage
\section{线性空间}
\subsection{定义}
\begin{formal}
\begin{definition}[线性空间]\label{def:线性空间}
    设有数域$\mathbb{K}$,$V$是一个非空集合
    ,在$V$上定义有加法运算：
    \begin{align*}
        + :
        V \times V & \longrightarrow V\\
        \left(\bm{\alpha},\bm{\beta}
        \right) &\longmapsto \bm{\alpha}+
        \bm{\beta}
    \end{align*}
    ,$\mathbb{K}$关于$V$定义有数乘运算
    \begin{align*}
        \cdot :
        \mathbb{K} \times V 
        &
        \longrightarrow V \\
        \left(k,\bm{\alpha}\right)
        &\longmapsto k\bm{\alpha}
    \end{align*}
    且满足下面的$8$条性质
    \begin{enumerate}[label=\textup{(\arabic*)}]
        \item 加法交换律：$\bm{\alpha}+\bm{\beta}
        =\bm{\beta}+\bm{\alpha}
        $
        \item 加法结合律：$
        \left(\bm{\alpha}+
        \bm{\beta}\right)+\bm{\gamma}
        = \bm{\alpha}+
        \left(\bm{\beta}+\bm{\gamma}\right)
    $
        \item 存在零向量$\bm{0}
        \in V$,使得$
        \forall \bm{\alpha} \in V$,
        均有$\bm{0}+\bm{\alpha}
        =\bm{\alpha}$
        \item $\forall \bm{\alpha}\in V
        $,存在其负向量$\bm{\beta}\in V$
            ,使得
        $\bm{\alpha} + \bm{\beta} = \bm{0}
        $
        \item 存在数乘的单位元$1\in\mathbb{K}
        $,使得$\forall\bm{\alpha}\in V
        $,均有$1\cdot\bm{\alpha} = 
        \bm{\alpha}$
        \item 数乘对向量的分配律：$k
        \left(\bm{\alpha}+\bm{\beta}\right)=
        k\bm{\alpha}+k\bm{\beta}
        $
        \item 数乘对数的分配律：$
        \left(k+l\right)\bm{\alpha}=
        k\bm{\alpha}+l\bm{\alpha}
        $
        \item 数乘的结合律：$
        k\left(l\bm{\alpha}\right)
        =
        \left(kl\right)\bm{\alpha}
        $
    \end{enumerate}
    则$V$称为数域$\mathbb{K}$上的线性空间
    或向量空间,记作$V_{\mathbb{K}}$.
\end{definition}
\end{formal}
\begin{brown}
\begin{example}
    \cref{def:n维行（列）向量空间}中的$\mathbb{K}^n$、$\mathbb{K}_n$
    均为$\mathbb{K}$上的的线性空间.这种线性空间具有普遍的代表性.
\end{example}
\end{brown}
\newpage
\begin{brown}
\begin{example}
    $\mathbb{K}$上的一元多项式全体
    \[
    \mathbb{K}\left[x\right]
    = \left\{a_nx^n+\cdots+
    a_1x+a_0\Bigg|\begin{matrix}
    x\text{为未定元}\qquad\quad\\
    a_i \in \mathbb{K},0\leqslant 
    i\leqslant n
    \end{matrix}\right\}
    \]
    是$\mathbb{K}$上的线性空间.其中运算按照寻常的运算定义
    ,即两个多项式的加法定义为同次项系数相加,一个数与一个多项式的数
    乘定义为
    这个数乘以多项式的每个系数.
    
    特别地,取
    其中次数不高于$n$的多项式全体$\mathbb{K}\left[x\right]$也是
    $\mathbb{K}$上的线性空间.
\end{example}
\end{brown}
\begin{brown}
\begin{example}
    设$C\left[0,1\right]$为$\left[0,1
    \right]$上全体连续函数的集合,定义
    \[
    \left(f+g\right)\left(x\right)=
    f\left(x\right)+g\left(x\right)
    \quad
    kf\left(x\right) = 
    \left(kf\right)\left(x\right)
    \]
    容易证明$C\left[0,1\right]$是
    $\mathbb{R}$的线性空间.
\end{example}
\end{brown}
\begin{brown}
\begin{example}\label{ex:矩阵全体}
    $M_{m \times n}\left(
        \mathbb{K}
    \right)$为$\mathbb{K}$上所有$m
    \times n$矩阵全体,其在矩阵的加法和数乘下是
    $\mathbb{K}$上的线性空间.
\end{example}
\end{brown}
\begin{red}
\begin{remark}
    \cref{ex:矩阵全体}$M_{n}\left(\mathbb{K}\right)$不仅是$\mathbb{K}$上的线性空间,还是一个$\mathbb{K}$-代数（\cref{def:代数}）.
\end{remark}
\end{red}
\begin{brown}
\begin{example}
    任意的$\mathbb{K}_1 \subseteq
    \mathbb{K}_2$,均有$\mathbb{K}_2$
    是$\mathbb{K}_1$的线性空间,其中
    \begin{align*}
        & + : a + b \quad a,b \in \mathbb{K}_2
        \\
        & \cdot : k \cdot a \quad 
        k \in \mathbb{K}_1,a \in 
        \mathbb{K}_2
    \end{align*}
    特别地,$\mathbb{K}_1
    =\mathbb{K}_2 =
    \mathbb{K}$时亦成立.
\end{example}
\end{brown}
\subsection{基本性质}
\begin{green}
\begin{corollary}[零向量唯一性]\label{cor:零向量唯一性}
    零向量唯一.
\end{corollary}
\begin{Proof}
设线性空间$V$的两个零向量$\bm{0}_1$、$\bm{0}_2$,
则\[
\bm{0}_1=\bm{0}_1+\bm{0}_2=\bm{0}_2
\]
唯一性得证.
\end{Proof}
\end{green}
\begin{green}
\begin{corollary}[负向量唯一性]\label{cor:负向量唯一性}
    负向量唯一.
\end{corollary}
\begin{Proof}
设$\bm{\alpha}\in V$,$\bm{\beta}_1$、$\bm{\beta}_2\in V$为
$\bm{\alpha}$的负向量.则
\begin{align*}
    \bm{\beta}_1&=\bm{\beta}_1+\bm{0}=\bm{\alpha}+\bm{\beta}_1+\bm{\beta}_2\\
    &=\bm{0}+\bm{\beta}_2=\bm{\beta}_2
\end{align*}
唯一性得证.
\end{Proof}
\end{green}
\begin{formal}
\begin{proposition}[线性空间的基本性质]\label{prop:线性空间的基本性质}
设$\mathbb{K}$上的线性空间$V$,
$\forall \bm{\alpha},
\bm{\beta},\bm{\gamma} 
\in V$,其除了\cref{def:线性空间}中的性质外还满足
\begin{enumerate}[label=\textup{(\arabic*)}]
    \item \cref{cor:零向量唯一性}零向量唯一
    \item \cref{cor:负向量唯一性}负向量唯一
    \item 加法消去律：$
    \bm{\alpha} + \bm{\beta}
    = \bm{\alpha}+\bm{\gamma}
    \Longrightarrow 
    \bm{\beta} = \bm{\gamma}
    $
    \item 数零与向量零：$0 \cdot \bm{\alpha}
    = \bm{0}
    $
    \item 零元：$k \cdot \bm{0} = \bm{0}$
    \item 负元；$ \left(-1\right)\bm{\alpha}
    = -\bm{\alpha}
    $
    \item 整性：$k\cdot \bm{\alpha} = \bm{0}
    \Longrightarrow k = 0$或$
    \bm{\alpha}=\bm{0}$
\end{enumerate}
\end{proposition}
\end{formal}
\newpage
\section{向量的线性关系}
\subsection{线性表示}
\begin{formal}
\begin{definition}[线性表示]\label{def:线性表示}
    设$\mathbb{K}$上的线性空间$V$,
    $\bm{\alpha}_1,\bm{\alpha}_2,\cdots,
    \bm{\alpha}_n \in 
    V_{\mathbb{K}}$,如果存在$
    k_1,k_2,\cdots,k_n \in \mathbb{K}$,
    使得
    \[
    \bm{\beta} = 
    k_1\bm{\alpha}_1+\cdots+
    k_n\bm{\alpha}_n
    \]
    称$\bm{\beta}$是$\bm{\alpha}_1,
    \bm{\alpha}_2,\cdots,\bm{\alpha}_n$的线性组合,
    或者说$\bm{\beta}$可以用$\bm{\alpha}_1
    ,\bm{\alpha}_2,\cdots,\bm{\alpha}_n$线性表示/表出.
    记作
    \[
    \bm{\beta}\hookrightarrow \left\{\bm{\alpha}_1,\bm{\alpha}_2,\cdots,\bm{\alpha}_n\right\}
    \]
\end{definition}
\end{formal}
\begin{brown}
\begin{example}
    $\mathbb{K}$上的$n$维标准单位行向量
    \[
    \bm{e}_i = \left(0,\cdots,1,\cdots,0\right)
    \quad
    \left(1\leqslant i\leqslant n\right)
    \]
    （其中$1$位于第$i$位）构成的向量组$\left\{\bm{e}_1,\bm{e}_2,\cdots
    ,\bm{e}_n\right\}$称为$\mathbb{K}$
    上的标准单位行向量组.显然
    任一$n$维行向量$\bm{\alpha}=
    \left(a_1,a_2,\cdots,a_n\right)$都可用其表示
    \[
    \bm{\alpha}=a_1\bm{e}_1+a_2\bm{e}_2+\cdots+a_n\bm{e}_n
    \]
\end{example}\end{brown}
\subsection{线性相关与线性无关}
\begin{formal}
\begin{definition}[线性相关与线性无关]\label{def:线性相关与线性无关}
    设$\mathbb{K}$上的线性空间$V$,
    $\bm{\alpha}_1,\bm{\alpha}_2,\cdots,
    \bm{\alpha}_n \in 
    V_{\mathbb{K}}$,
    若存在不全为零的$k_1,k_2,\cdots,k_n \in 
    \mathbb{K}$使得\[
    k_1\bm{\alpha}_1+k_2\bm{\alpha}_2+\cdots
    +k_n\bm{\alpha}_n = \bm{0}
    \]
    称$\bm{\alpha}_1,\bm{\alpha}_2,\cdots,
    \bm{\alpha}_n $是线性相关的；
    反之若不存在这样的一组
    $k_1,k_2,\cdots,k_n\in\mathbb{K}$,则称
    $\bm{\alpha}_1,\bm{\alpha}_2,\cdots,
    \bm{\alpha}_n $是线性无关的.
\end{definition}\end{formal}
\begin{formal}
\begin{definition}[线性无关的等价定义]\label{def:线性无关的等价定义}
    \cref{def:线性相关与线性无关}线性无关等价定义为：若存在
    $k_1,k_2,\cdots,k_n\in V$使得
    \[
    k_1\bm{\alpha}_1 +k_2\bm{\alpha}_2+\cdots +
    k_n\bm{\alpha}_n = \bm{0}
    \]
    那么$k_i = 0\left(i = 1,2,\cdots,n
    \right).$
\end{definition}\end{formal}
\begin{red}
\begin{remark}
    容易得到,方程$\bm{A}\bm{x} = \bm{\beta}$
有非零解的充要条件是
\[
\bm{A} = \left(\bm{\alpha}_1,\cdots
,\bm{\alpha}_n\right)
\]
中$\bm{\alpha}_1,\cdots,\bm{\alpha}_n$
线性相关.
\end{remark}
\end{red}
\begin{formal}
\begin{theorem}[线性关系的非唯一性]\label{thm:线性关系的非唯一性}
    线性相关或线性无关依赖于基域$\mathbb{K}
    $的选定$\left(k_i \in \mathbb{K}
    \right)$.
\end{theorem}
\end{formal}
\begin{brown}
\begin{example}
    $\mathbb{R} \subseteq
    \mathbb{C}$,易证$\left\{1,\mathrm{i}
    \right\}$在$\mathbb{R}$线性无关,
    在$\mathbb{C}$线性相关.
\end{example}\end{brown}
\begin{green}
\begin{corollary}[含有零向量的向量组]\label{cor:含有零向量的向量组}
    包含$\bm{0}$的向量组一定线性相关.
\end{corollary}
\end{green}
\begin{green}
\begin{corollary}[仅一个向量的向量组]\label{cor:仅一个向量的向量组}
    \[
    A = \left\{\bm{\alpha}\right\} \text{
        线性无关
    } \Longleftrightarrow 
    \bm{\alpha} \neq \bm{0}
    \]
\end{corollary}
\end{green}
\begin{formal}
\begin{theorem}[向量组的大小关系]\label{thm:向量组的大小关系}
    $n > m > s$,如果$\left\{\bm{\alpha}_1
    ,\bm{\alpha}_2,\cdots,\bm{\alpha}_m\right\}$线
    性相关,那么$\left\{\bm{\alpha}_1,
    \cdots,\bm{\alpha}_m,\cdots,
    \bm{\alpha}_n\right\}$也线性相关；
    如果$\left\{\bm{\alpha}_1,\bm{\alpha}_2
    ,\cdots,\bm{\alpha}_m\right\}$,线
    性无关,那么$\left\{\bm{\alpha}_1,
    \cdots,\bm{\alpha}_s\right\}$也
    线性无关.
\end{theorem}
\end{formal}
\begin{green}
\begin{corollary}[余向量表示]\label{cor:余向量表示}
    $\bm{\alpha}_1,\bm{\alpha}_2,\cdots,
    \bm{\alpha}_n$线性相关等价于
    存在$\bm{\alpha}_i\left(
        1 \leqslant i
        \leqslant n
    \right)$是$\left\{
        \bm{\alpha}_1,\cdots
        ,\widehat{\bm{\alpha}_i},
        \cdots,\bm{\alpha}_n
    \right\}$的线性组合.
\end{corollary}
\end{green}
\begin{green}
\begin{corollary}[加入向量的线性关系]\label{cor:加入向量的线性关系}
    设$\bm{\alpha}_1,\bm{\alpha}_2,\cdots,
    \bm{\alpha}_n,\bm{\beta}
    \in  V_{\mathbb{K}}$
    ,且$\bm{\alpha}_1,\bm{\alpha}_2,\cdots
    ,\bm{\alpha}_n$是线性无关的,
    那么一定有
    要么$\bm{\alpha}_1,\bm{\alpha}_2,\cdots,
    \bm{\alpha}_n,\bm{\beta}
    $线性无关,要么$\bm{\beta}$
    是$\bm{\alpha}_1,\bm{\alpha}_2,\cdots,
    \bm{\alpha}_n$的线性组合.
\end{corollary}
\begin{Proof}
考虑线性相关与线性无关的相补即可.
\end{Proof}
\end{green}
\begin{formal}
\begin{theorem}[线性表示唯一]\label{thm:线性表示唯一}
    已知$\bm{\alpha}_1,\bm{\alpha}_2,\cdots,
    \bm{\alpha}_n,\bm{\beta}
    \in  V_{\mathbb{K}}$且
    \[
    \bm{\beta} = k_1\bm{\alpha}_1+k_2\bm{\alpha}_2
    +\cdots+k_n\bm{\alpha}_n
    \]
    则该表示形式唯一当且仅当$\bm{\alpha}_1,\bm{\alpha}_2,\cdots,\bm{\alpha}_n$线性无关.
\end{theorem}\end{formal}
\begin{formal}
\begin{theorem}[线性组合的传递性]\label{thm:线性组合的传递性}
    设有向量组$A = \left\{
        \bm{\alpha}_1,\bm{\alpha}_2,\cdots,
        \bm{\alpha}_m
    \right\}$、$
    B = \left\{
        \bm{\beta}_1,\bm{\beta}_2,\cdots,
        \bm{\beta}_n
    \right\}$、$C = 
    \left\{\bm{\gamma}_1,\bm{\gamma}_2,\cdots
    ,\bm{\gamma}_p\right\}$
    ,若$B$中的向量均为$A$中的向量
    的线性组合且$C$中的向量均为
    $B$中的向量的线性组合,那么
    $C$中的向量均可用$A$中的向量
    线性表示.即若
    \[
    C\hookrightarrow B \hookrightarrow A
    \]
    则
    \[
    C\hookrightarrow A
    \]
\end{theorem}
\end{formal}
\begin{brown}
\begin{example}
    $\bm{\alpha} = \left(a_1,a_2,
    \cdots,a_n\right)$与$
    \bm{\beta} = \left(b_1,b_2,\cdots,
    b_n\right)$线性相关的充分必要条件是
    $a_i = c b_i(i = 1,2,
    \cdots,n,c \in\mathbb{R})$成立.
\end{example}
\end{brown}
\begin{brown}
\begin{example}
    $V = \mathbb{R}^2$中
    $\overrightarrow{OA} =\left(
        a_1,a_2
    \right)$、$\overrightarrow{OB}=
    \left(b_1,b_2\right)
     $线性相关等价于$O$、$A$、$B$共线；
     $\overrightarrow{OA}
     =\left(a_1,a_2\right) $、
     $\overrightarrow{OB} = \left(b_1
     ,b_2\right)
     $线性无关等价于
     $\bigtriangleup OAB$非退化
     ,这时
     \[
     S_{\bigtriangleup OAB} =
     \frac{\left|\overrightarrow{OA}
     \times \overrightarrow{OB} \right|
     }{2} = \frac{
        \begin{vmatrix}
            a_1 & a_2 \\
            b_1 & b_2
        \end{vmatrix}
     }{2}
     \neq 0
     \]

     $V = \mathbb{R}^3$中
    $\overrightarrow{OA} =\left(
        a_1,a_2,a_3
    \right)$、$\overrightarrow{OB}=
    \left(b_1,b_2,b_3\right)
     $、
     $\overrightarrow{OC} = 
     \left(c_1,c_2,c_3\right) $线
     性相关等价于$O$、$A$、$B$、$C$共面
     ；
     $\overrightarrow{OA}
     =\left(a_1,a_2,a_3\right) $、
     $\overrightarrow{OB} = \left(b_1
     ,b_2,b_3\right)
     $、$
     \overrightarrow{OC} = 
     \left(c_1,c_2,c_3\right) $线
     性无关等价于四面体$ OABC$非退化
     （或者说该平行六面体非退化）
     ,这时
     \[
        \begin{vmatrix}
            a_1 & a_2 & a_3 \\
            b_1 & b_2 & b_3 \\
            c_1 & c_2 & c_3
        \end{vmatrix}
     \neq 0
     \]
\end{example}\end{brown}
\begin{green}
\begin{corollary}[向量组线性关系的行列式判断]\label{cor:向量组线性关系的行列式判断}
    $\bm{\alpha}_1 = \left(a_{11},a_{21},
    \cdots,a_{n1}\right),\bm{\alpha}_2=
    \left(a_{12},a_{22},\cdots,a_{2n}\right),\cdots,\bm{\alpha}_n
    =\left(a_{1n},a_{2n},\cdots,a_{nn}\right)
    \in \mathbb{K}^n$线性无关等价于
    \[
    \left|\bm{\alpha}_1,\bm{\alpha}_2,\cdots
    ,\bm{\alpha}_n\right|=
    \begin{vmatrix}
        a_{11} & a_{12} & \cdots & a_{1n} \\
        a_{21} & a_{22} & \cdots & a_{2n} \\
        \vdots & \vdots & s & \vdots \\
        a_{n1} & a_{n2} & \cdots & a_{nn}
    \end{vmatrix}
    \neq 0
    \]
    即是矩阵可逆,也就是满秩.反之则是缺秩,对应线性相关.
\end{corollary}
\end{green}
\newpage
\section{向量组的秩}
\subsection{向量族与向量组}
\begin{formal}
\begin{definition}[向量族与向量组]\label{def:向量族与向量组}
    在向量空间$V_{\mathbb{K}}$中,向量族
    是$V$中的向量的集合（可能有无穷多个
    且允许存在重复的元）,向量组是
    $V$中有限个向量的集合.

    向量族或向量组$S$包含的向量个数一般记作$\#S$.
\end{definition}
\end{formal}
\begin{formal}
\begin{definition}[极大线性无关组]\label{def:极大线性无关组}
    设$S$为$V_{\mathbb{K}}$中一个向量族,
    若其中存在一组向量$
    \left\{
        \bm{\alpha}_1,\bm{\alpha}_2,\cdots,
        \bm{\alpha}_r
    \right\}$使得
    \begin{enumerate}[label=\textup{(\arabic*)}]
        \item $\bm{\alpha}_1,\bm{\alpha}_2,\cdots,\bm{\alpha}_r
        $线性无关
        \item $\forall \bm{\beta} \in S,
        \bm{\beta} = k_1\bm{\alpha}_1 +k_2\bm{\alpha}_2+
        \cdots +k_r\bm{\alpha}_r
        $
    \end{enumerate}
    则称向量组$\left\{
        \bm{\alpha}_1,\bm{\alpha}_2,\cdots,
        \bm{\alpha}_r
    \right\}$为$S$的一个极大（线性）无关组.
\end{definition}
\end{formal}
\begin{green}
\begin{corollary}[极大无关组的存在性]\label{cor:极大无关组的存在性}
    包含非零向量的向量组必有极大无关组.
\end{corollary}
\begin{Proof}
对$S$包含的向量个数$k$作归纳.

$k=1$时,设$S=\left\{\bm{\alpha}\right\}$且$
\bm{\alpha}\neq \bm{0}$,那么$\left\{
    \bm{\alpha}\right\}$就是极大无关组.
    
    一般地,若$S$中$k$个向量时线性无关的,
    那么它们即是极大无关组.若是线性相关的即至少存在
    一个向量可用其余向量线性表出,不妨设$\bm{\alpha}$
    可用$S\backslash\left\{\bm{\alpha}\right\}$（含有
    $k-1$个向量且至少一个非零向量）线性表出.
    于是$S\backslash \left\{\bm{\alpha}\right\}$存在极大无关组
    $\left\{\bm{\alpha}_1,
    \bm{\alpha}_2,\cdots,\bm{\alpha}_r\right\}$,那么
    由线性组合的传递性知$\bm{\alpha}$也可以用向量
    $\bm{\alpha}_1,\bm{\alpha}_2,\cdots,\bm{\alpha}_r$
    线性表出,于是$\left\{\bm{\alpha}_1,\bm{\alpha}_2,\cdots,\bm{\alpha}_r\right\}$
    就是$S$的极大无关组.
\end{Proof}
\end{green}
\begin{green}
\begin{lemma}[多少之间的线性关系]\label{lem:多少之间的线性关系}
    $A$、$B$两个向量组分别包含$r$、$s$
    个向量,且$A$中的向量均可由$B$中的向量
    表示,若$A$中向量线性无关则必有$r
    \leqslant s$即\[\#A \leqslant \#B\]
\end{lemma}
\begin{Proof}
设\begin{align*}
    A=\left\{\bm{\alpha}_1,\bm{\alpha}_2,\cdots
    ,\bm{\alpha}_r\right\},\#A=r\\
    B=\left\{\bm{\beta}_1,\bm{\beta}_2,\cdots,\bm{\beta}_s\right\},\#B=s
\end{align*}
考虑反证法,设$r>s$.设
\[
\bm{\alpha}_1=\lambda_1\bm{\beta}_1+\lambda_2\bm{\beta}_2+\cdots+\lambda_s\bm{\beta}_s
\]
因为$A$线性无关则$\bm{\alpha}_1\neq \bm{0}$,于是
$\lambda_1,\lambda_2,\cdots,\lambda_s$不全为零.不妨设
$\lambda_1\neq 0$,则
\[
\bm{\beta}_1=\frac{1}{\lambda_1}\bm{\alpha}_1-
\frac{\lambda_2}{\lambda_1}\bm{\beta}_2-\cdots-
\frac{\lambda_s}{\lambda_1}\bm{\beta}_s
\]

因为$A\backslash \left\{\bm{\alpha}_1\right\}$中的向量
均可用$B$线性表示,且由上知$B$中的
$\bm{\beta}_1$可用
$\left\{\bm{\alpha}_1,\bm{\beta}_2,\cdots,\bm{\beta}_s\right\}$
线性表示.

于是$\forall 1<i\leqslant r$,$\bm{\alpha}_i$是
$\bm{\alpha}_1,\bm{\beta}_2,\cdots,\bm{\beta}_s$的线性组合.

考虑归纳法,设$\forall k < i \leqslant r$,
$\bm{\alpha}_i$是
$\left\{\bm{\alpha}_1,
\bm{\alpha}_2,\cdots,\bm{\alpha}_k,
\bm{\beta}_{k+1},\cdots,\bm{\beta}_s\right\}$的线性组合.
不妨取
\[
\bm{\alpha}_{k+1}=\mu_{1}\bm{\alpha}_1+\mu_2\bm{\alpha}_2+\cdots+\mu_k\bm{\alpha}_k+
\mu_{k+1}\bm{\beta}_{k+1}+\cdots+\mu_s\bm{\beta}_s
\]
若$\mu_{k+1}=\cdots=\mu_{s}=0$则
$\bm{\alpha}_{k+1}$是
$\bm{\alpha}_1,\bm{\alpha}_2,
\cdots,\bm{\alpha}_k$的线性组合,与$A$线性无关
矛盾,于是
$\mu_{k+1},\cdots,\mu_s$不全为零.不妨设$\mu_{k+1}\neq 0$,则
\[
\bm{\beta}_{k+1}=
-\frac{\mu_1}{\mu_{k+1}}\bm{\alpha}_1-
\frac{\mu_2}{\mu_{k+1}}\bm{\alpha}_2-\cdots
-\frac{\mu_k}{\mu_{k+1}}\bm{\alpha}_k+
\frac{1}{\mu_{k+1}}\bm{\alpha}_{k+1}-
\frac{\mu_{k+2}}{\mu_{k+1}}\bm{\beta}_{k+2}-
\cdots-\frac{\mu_{s}}{\mu_{k+1}}\bm{\beta}_s
\]
那么
\begin{align*}
    \left\{\bm{\alpha}_{k+2},\cdots,\bm{\alpha}_r\right\}
    &\hookrightarrow 
    \left\{\bm{\alpha}_1,
\bm{\alpha}_2,\cdots,\bm{\alpha}_k,
\bm{\beta}_{k+1},\cdots,\bm{\beta}_s\right\}\\
&\hookrightarrow
\left\{\bm{\alpha}_1,
\bm{\alpha}_2,\cdots,\bm{\alpha}_k,\bm{\alpha}_{k+1},
\bm{\beta}_{k+2},\cdots,\bm{\beta}_s\right\}
\end{align*}
于是$\forall s < i \leqslant r$,
$\bm{\alpha}_i$是$\left\{\bm{\alpha}_1,\bm{\alpha}_2,\cdots,\bm{\alpha}_s\right\}$
的线性组合,与$A$线性无关矛盾.于是证毕.
\end{Proof}
\end{green}
\begin{green}
\begin{corollary}[逆否命题]\label{cor:逆否命题}
\cref{lem:多少之间的线性关系}的逆否命题为若多的向量组可以用少的
向量组线性表示,则多的向量组一定线性相关.
\end{corollary}
\end{green}
\begin{green}
\begin{lemma}[相互表出的线性无关向量组的向量个数关系]\label{lem:相互表出的线性无关向量组的向量个数关系}
    设$A$、$B$为线性无关的向量组,
    且二者包含的向量可以互相表示,那么必有
    二者所含向量个数相同即$\#A=\#B$.
\end{lemma}
\begin{Proof}
由\cref{lem:多少之间的线性关系}得到.
\end{Proof}
\end{green}
\begin{green}
\begin{corollary}[极大无关组向量个数唯一]\label{cor:极大无关组向量个数唯一}
    一个向量族极大线性无关组中向量的选取不一定是唯一的,但其中的向量的个数一定是唯一的.
\end{corollary}
\begin{Proof}
由\cref{lem:多少之间的线性关系}和\cref{lem:相互表出的线性无关向量组的向量个数关系}得到.
\end{Proof}
\end{green}
\subsection{向量组的秩}
\begin{formal}
\begin{theorem}[极大无关组向量个数唯一]\label{thm:极大无关组向量个数唯一}
    $A$、$B$均为$S$的极大无关组,则$A$
    与$B$所包含的向量个数相同.
\end{theorem}
\begin{Proof}
即\cref{cor:极大无关组向量个数唯一}.
\end{Proof}
\end{formal}
\begin{formal}
\begin{definition}[秩]\label{def:秩}
    向量族$S$的极大无关组
    包含的向量个数
    称为$S$的秩,记作$\mathrm{r}\left(S\right)$
    或者$\mathrm{rank}\left(S\right)$.
\end{definition}
\end{formal}
\begin{red}
    \begin{remark}
        由\cref{cor:极大无关组向量个数唯一},秩的值与极大无关组的选取无关.
    \end{remark}
\end{red}
\begin{red}
\begin{remark}
    特别地,约定所有元为$\bm{0}$的向量族的秩为$0$.
\end{remark}
\end{red}
\begin{formal}
\begin{definition}[等价]\label{def:等价}
    可以互相线性表示的向量组称为等价.
\end{definition}
\end{formal}
\begin{formal}
\begin{theorem}[同秩]\label{thm:同秩}
    等价的向量组的秩相同,称为同秩.
\end{theorem}
\begin{Proof}
如果向量组$A$、$B$全是零向量,明显成立.
若二者均至少有一个非零向量,设它们的极大无关组分别为
$A_1$、$B_1$,则
$\mathrm{rank}\left(A\right)=\#A_1,\mathrm{rank}\left(B\right)=\#B_1$.
因为$A_1$、$B_1$线性无关,且有
\begin{align*}
    A_1\subseteq A \hookrightarrow B \hookrightarrow B_1\\
    B_1\subseteq B \hookrightarrow A \hookrightarrow A_1
\end{align*}
于是$\#A_1=\#B_1$即$\mathrm{rank}\left(A\right)=\mathrm{rank}\left(B\right).$
\end{Proof}
\end{formal}
\subsection{基}
\begin{red}
    \begin{remark}
        如果$S = V_{\mathbb{K}}$,那么
\[
\begin{cases*}
    \text{极大无关组——线性空间的基}
    \\
    \text{极大无关组的秩——线性空间的维数}
\end{cases*}
\]
    \end{remark}
\end{red}
\begin{formal}
\begin{definition}[基]\label{def:基}
    设$V$是$\mathbb{K}$上的线性空间,
    若其中存在线性无关的一组向量
    $\left\{\bm{e}_1,\bm{e}_2
    ,\cdots,\bm{e}_n\right\}$
    使得$V$中任一向量均可表示为
    这组向量的线性组合,则称$\left\{
        \bm{e}_1,\bm{e}_2,\cdots
        ,\bm{e}_n
    \right\}$为$V$的一组基,$V$
    称为$n$维线性空间（具有维数$n$）,
    记作
    \[
    \dim_{\mathbb{K}}\,V = n
    \]
    这组基记作
    \[
    \left\{
        \bm{e}_1,\cdots
        ,\bm{e}_n
    \right\}=
    V /
    \mathbb{K}
    \]
    若不存在这样的有限个向量组成的一组
    基,那么称$V$是无限维线性空间.
\end{definition}
\end{formal}
\begin{red}
\begin{remark}
    需要说明的是,虽然我们主要研究有限维线性空间,但
无限维线性空间仍然具有基,
只是需要修改线性表示和线性无关的定义并使用
选择公理或Zorn引理.
\end{remark}
\end{red}
\begin{brown}
\begin{example}
    \cref{def:n维行（列）向量空间}中的
    \[
    \dim_{\mathbb{K}}\,\mathbb{K}^n
    =\dim_{\mathbb{K}}\,\mathbb{K}_n
    = n
    \]
\end{example}
\end{brown}
\begin{brown}
\begin{example}
    \[
    \left\{1,\mathrm{i} = \sqrt{-1}\right\} = 
    \mathbb{C} /
    \mathbb{R} 
    \Longrightarrow 
    \dim_{\mathbb{R}}\,
    \mathbb{C} = 2
    \]
\end{example}\end{brown}
\begin{green}
\begin{corollary}[超过维数的向量组]\label{cor:超过维数的向量组}
    $n$维空间$V_{\mathbb{K}}$中
    任一包含向量个数多于
    $n$的向量组一定线性相关.
\end{corollary}
\end{green}
\begin{formal}
\begin{theorem}[基的判定定理]\label{thm:基的判定定理}
    设$\bm{e}_1,\bm{e}_2,\cdots,\bm{e}_n $是$\mathbb{K}$上的
    $n$维线性空间$V$中的$n$个向量,若
    适合下面之一则$\left\{
        \bm{e}_1,\bm{e}_2,\cdots,\bm{e}_n
    \right\}$是
    $V$的一组基
    \begin{enumerate}[label=\textup{(\arabic*)}]
       \item $\bm{e}_1,\bm{e}_2,\cdots,\bm{e}_n$
        线性无关
       \item $V$中任一向量均可表示为这组向量
       $\left\{
        \bm{e}_1,\bm{e}_2,\cdots,\bm{e}_n
    \right\}$
       的线性组合
    \end{enumerate}   
\end{theorem}
\begin{Proof}
$(1)$$\forall \bm{\alpha}\in V$,$\bm{e}_1,\bm{e}_2,\cdots,\bm{e}_n,\bm{\alpha}$线性相关,
故$\bm{\alpha}$是$\bm{e}_1,\bm{e}_2,\cdots,\bm{e}_n$的线性组合,证毕.

$(2)$设$\left\{\bm{\alpha}_1,\bm{\alpha}_2,\cdots
,\bm{\alpha}_r\right\}$为$\left\{
    \bm{e}_1,\bm{e}_2,\cdots,\bm{e}_n
\right\}$的极大无关组.即
\[
V\hookrightarrow \left\{\bm{e}_1,\bm{e}_2,\cdots
,\bm{e}_n\right\}\hookrightarrow
\left\{\bm{\alpha}_1,\bm{\alpha}_2,\cdots
,\bm{\alpha}_r\right\}
\]
于是$\left\{\bm{\alpha}_1,\bm{\alpha}_2,\cdots
,\bm{\alpha}_r\right\}$是$V$的一组基
,$\dim V=r=n$.证毕.
\end{Proof}
\end{formal}
\begin{formal}
\begin{theorem}[基扩张定理]\label{thm:基扩张定理}
    设数域$\mathbb{K}$上的
    $n$维线性空间$V_{\mathbb{K}}$
    中的$m\left(m < n\right)$
    个线性无关的向量为
    $\bm{v}_1,\bm{v}_2
    ,\cdots,\bm{v}_m$,$\left\{
        \bm{e}_1,\bm{e}_2,\cdots
        ,\bm{e}_n
    \right\}$是
    该空间的一组基,则必定可以
    从这组基
    再选出$n - m$个向量来共同组成
    一组基
    \[
    \left\{\bm{v}_1,\bm{v}_2,
    \cdots,
    \bm{v}_m,\bm{e}_{i_1},\bm{e}_{i_2},
    \cdots,\bm{e}_{i_{n - m}}\right\}
    \]
\end{theorem}
\begin{Proof}
首先证明$\exists 1\leqslant i \leqslant n $使得
$\bm{v}_1,\bm{v}_2,\cdots,
\bm{v}_m,\bm{e}_i$线性无关.

考虑反证法,设$\forall 1
\leqslant i \leqslant n $,
$\bm{v}_1,\bm{v}_2,\cdots,
\bm{v}_m,\bm{e}_i$线性相关,于是$\bm{e}_i$是
$\bm{v}_1,\bm{v}_2,\cdots,
\bm{v}_m$的线性组合.即
\[
\left\{\bm{e}_1,\bm{e}_2,\cdots,\bm{e}_n\right\}
\hookrightarrow \left\{
    \bm{v}_1,\bm{v}_2,\cdots,\bm{v}_m
\right\}
\]
且$\left\{\bm{e}_1,\bm{e}_2,\cdots,\bm{e}_n\right\}
$是线性无关的,于是$n \leqslant m$,矛盾.得证.

不妨设$\bm{v}_1,\bm{v}_2,\cdots,\bm{v}_m,\bm{e}_1$是线性无关的.
\begin{enumerate}[label=\textup{(\arabic*)}]
\item $m+1=n$时,这成为一组基.

\item $m+1<n$时,仿照上面不断做下去即可.
\end{enumerate}
证毕.
\end{Proof}
\end{formal}
\begin{formal}
\begin{theorem}[基扩张定理]\label{thm:基扩张定理2}\cref{thm:基扩张定理}基扩张定理的等价表述为
    \begin{enumerate}[label=\textup{(\arabic*)}]
        \item $n$维线性空间$V_{\mathbb{K}}$中任意$m\left(m < n\right)$个线性无关的向量均可扩张为$V$的一组基
        \item $V$的任一子空间的基
    均可扩张为$V$的一组基.
    \end{enumerate}
\end{theorem}
\end{formal}
\newpage
\section{矩阵的秩}
\subsection{定义}
\begin{formal}
\begin{definition}[行秩与列秩]\label{def:行秩与列秩}
    设矩阵$\bm{A}_{m \times n}$,
    有行分块\[
    \bm{A}=\begin{pmatrix}
        \bm{\alpha}_1\\
        \bm{\alpha}_2\\
        \vdots\\
        \bm{\alpha}_m
    \end{pmatrix}
    \]
    和列分块
    \[
    \bm{A}=\left(
        \bm{\beta}_1,\bm{\beta}_2,\cdots,\bm{\beta}_n
    \right)
    \]
    那么这$m$个行向量组成的向量组的秩$\mathrm{rank}\left(
        \bm{\alpha}_1,\bm{\alpha}_2,\cdots,\bm{\alpha}_m
    \right)$称为
    矩阵$\bm{A}$的行秩,$n$个列向量
    组成的向量组的秩
    $\mathrm{rank}\left(
        \bm{\beta}_1,\bm{\beta}_2,\cdots,\bm{\beta}_n
    \right)$称为矩阵$\bm{A}$的列秩.
\end{definition}
\end{formal}
\begin{green}
\begin{lemma}[极大无关组的指标保持]\label{lem:极大无关组的指标保持}
    设矩阵$\bm{A}_{m \times n} = 
    \left(\bm{\beta}_1,\bm{\beta}_2,\cdots,
    \bm{\beta}_n\right)$,$\bm{Q}$为
    $m$阶非异阵,若$\left\{\bm{\beta}
    _{i_1},\cdots,
    \bm{\beta}_{i_r}\right\}$是$\bm{A}$
    的列向量的极大无关组,则$
    \left\{
        \bm{Q}\bm{\beta}_{i_1},\bm{Q\beta}_{i_2}
        ,\cdots,
        \bm{Q}\bm{\beta}_{i_r}
    \right\}$是$\bm{QA}
    = \left(\bm{Q}\bm{\beta}_1,\bm{Q\beta}_{2},\cdots
    ,\bm{Q}\bm{\beta}_n\right)$
    的列向量的极大无关组.

    即初等行变换
保持矩阵列向量极大无关组的列指标.
\end{lemma}
\begin{Proof}
参考\cref{thm:初等变换不改变秩}的证明.
\end{Proof}
\end{green}
\begin{formal}
\begin{theorem}[初等变换不改变秩]\label{thm:初等变换不改变秩}
    初等变换不改变矩阵的行秩和列秩.同时,分块初等变换
    也不改变分块矩阵的秩.
\end{theorem}
\begin{Proof}
下面证明列秩.

记$\mathrm{r}_c\left(\bm{A}\right)=
\mathrm{r}\left(\bm{\beta}_1,\bm{\beta}_2,
\cdots,\bm{\beta}_n\right)$为$\bm{A}$列秩.

先证明初等列变换,设初等矩阵$\bm{Q}$,因为
\begin{enumerate}[label=\textup{(\arabic*)}]
    \item $\bm{AP}_{ij}=\left(\bm{\beta}_1,\cdots,\bm{\beta}_j,\cdots,\bm{\beta}_i,\cdots,\bm{\beta}_n\right)$
    \item $\bm{AP}_i\left(c\right)=\left(\bm{\beta}_1,\cdots,c\bm{\beta}_i,\cdots,\bm{\beta}_n\right)\left(c \neq 0\right)$
    \item $\bm{AT}_{ji}\left(c\right)=\left(\bm{\beta}_1,\cdots,\bm{\beta}_i,\cdots,\bm{\beta}_j+c\bm{\beta}_i,\cdots,\bm{\beta}_n\right)$
\end{enumerate}

容易看出$\bm{AQ}$的列向量均为$\bm{A}$的列向量的线性组合.我们考虑
\[
\bm{A}=\left(\bm{AQ}\right)\bm{Q}^{-1}
\]
那么上面的论断反过来也是正确的.
从而$\bm{A}$和$\bm{AQ}$的列向量组一定
是等价的.于是
$\mathrm{r}_c\left(\bm{A}\right)=\mathrm{r}_c\left(\bm{AQ}\right)$,证毕.

然后证明初等行变换.我们证明更强的结论即\cref{lem:极大无关组的指标保持}.

首先证明$\bm{Q\beta}_{i_1},\bm{Q\beta}_{i_2},\cdots,\bm{Q\beta}_{i_r}$线性无关.
设
\[
\lambda_1\bm{Q\beta}_{i_1}+\lambda_2\bm{Q\beta}_{i_2}+\cdots+\lambda_r\bm{Q\beta}_{i_r}=\bm{0}
\]
即
\[
\bm{Q}\left(\lambda_1\bm{\beta}_{i_1}+\lambda_2\bm{\beta}_{i_2}+\cdots+\lambda_r\bm{\beta}_{i_r}\right)=\bm{0}
\]
此处$\bm{Q}$非异,可约
\[
\lambda_1\bm{\beta}_{i_1}+\lambda_2\bm{\beta}_{i_2}+\cdots+\lambda_r\bm{\beta}_{i_r}=\bm{0}
\]
于是$\lambda_1=\lambda_2=\cdots=\lambda_r=0$.于是结论成立.

再证$\bm{Q\beta}_j$均为$\bm{Q\beta}_{i_1},\bm{Q\beta}_{i_2},\cdots,\bm{Q\beta}_{i_r}$的
线性组合.因为$\bm{\beta}_j$是
$\bm{\beta}_{i_1},\bm{\beta}_{i_2},\cdots,\bm{\beta}_{i_r}$
的线性组合.于是得证.

\cref{lem:极大无关组的指标保持}证毕.

令$\bm{Q}$为初等矩阵,由\cref{lem:极大无关组的指标保持}立刻得证
$\mathrm{r}_c\left(\bm{A}\right)=
\mathrm{r}_c\left(\bm{QA}\right)$.注意此时考虑特例
$\bm{A}=\bm{O}$亦成立.

综上,证毕.
\end{Proof}
\end{formal}
\begin{formal}
\begin{theorem}[行列秩关系]
    矩阵的行秩等于列秩,称为矩阵的秩,记作
    $\mathrm{r}\left(\bm{A}\right)$或
    $\mathrm{rank}\left(\bm{A}\right)$
    .
\end{theorem}
\begin{Proof}
因为任一矩阵均有相抵标准型
\[
\begin{pmatrix}
    \bm{I}_r & \bm{O} \\
    \bm{O} & \bm{O} 
\end{pmatrix}
\]再结合\cref{thm:初等变换不改变秩},证毕.
\end{Proof}
\end{formal}
\begin{green}
\begin{corollary}[转置的秩]\label{cor:转置的秩}
    $\bm{A} \in 
    M_{m \times n}\left(
        \mathbb{K}
    \right)$,则$\mathrm{r}\left(
        \bm{A}\right) = \mathrm{r}
        \left(\bm{A}^{\prime}\right)$
\end{corollary}
\end{green}
\begin{green}
\begin{corollary}[与可逆阵积的秩]\label{cor:与可逆阵积的秩}
    $\bm{A} \in 
    M_{m \times n}\left(
        \mathbb{K}
    \right)$,$\bm{P}$、
    $\bm{Q}$分别为$m$、$n$阶可逆阵
    ,则\[
    \mathrm{r}\left(\bm{PAQ}
    \right) = \mathrm{r}\left(
        \bm{A}
    \right)
    \]
\end{corollary}
考虑到\cref{thm:初等变换与初等矩阵},再根据\cref{thm:初等变换不改变秩},证毕.
\end{green}
\begin{green}
    \begin{corollary}[相抵标准型]\label{cor:相抵标准型}
    对任意$\bm{A} \in 
    M_{m \times n}\left(
        \mathbb{K}
    \right)$,$\mathrm{r}\left(\bm{A}
    \right)= r$,总存在
    $m$、$n$阶可逆阵$\bm{P}
    \in M_m\left(\mathbb{K
    }\right)$、
    $\bm{Q}\in
    M_n\left(
        \mathbb{K}
    \right)$使得
    \[
    \bm{PAQ} = \begin{pmatrix}
        \bm{I}_r & \bm{O} \\
        \bm{O} & \bm{O}
    \end{pmatrix}
    \]
\end{corollary}
\begin{Proof}
由\cref{lem:相抵标准型}和\cref{cor:与可逆阵积的秩}得到.
\end{Proof}
\end{green}
\begin{green}
    \begin{corollary}[相抵与秩]\label{cor:相抵与秩}
    如果$\bm{A}$,$\bm{B}\in M_{
        m \times n
    }\left(\mathbb{K}\right)$,
    则$\bm{A} \sim \bm{B} \Longleftrightarrow 
    \mathrm{r}\left(\bm{A}\right)
    = \mathrm{r}\left(\bm{B}
    \right)$
\end{corollary}
\begin{Proof}
根据\cref{cor:与可逆阵积的秩}或\cref{cor:相抵标准型}得到.
\end{Proof}
\end{green}
\begin{formal}
    \begin{definition}[满秩]\label{def:满秩}
    如果矩阵所有的行向量线性无关,称为
    行满秩；如果矩阵所有列向量线性无关,
    称为列满秩.
\end{definition}
\end{formal}
\begin{formal}
\begin{theorem}[矩阵可逆与满秩]\label{thm:矩阵可逆与满秩}
特别地,对于$n$阶方阵$\bm{A}$
\[
\left|\bm{A}\right| \neq 0
\Longleftrightarrow 
\bm{A}\text{可逆}\Longleftrightarrow 
\mathrm{r}\left(\bm{A}\right)
= n\Longleftrightarrow 
\bm{A}\text{是满秩阵}\begin{cases*}
    \text{行满秩}\Longleftrightarrow
    \text{行向量线性无关} \\
    \text{列满秩}\Longleftrightarrow
    \text{列向量线性无关} 
\end{cases*}
\]
\end{theorem}
\begin{Proof}
我们证明方阵可逆与满秩的等价.

一方面,
$\mathrm{rank}\left(\bm{A}\right)=n$时$\bm{A}\sim \bm{I}_n$
,可逆.另一方面,$\bm{A}$可逆时
$\bm{A}=\bm{AI}_n\Longrightarrow
\mathrm{rank}\left(\bm{A}\right)=\mathrm{rank}\left(\bm{I}_n\right)
=n$

证毕.
\end{Proof}
\end{formal}
首先给出一个引理,其与\cref{thm:基的判定定理}基的判定定理相似.
\begin{green}
\begin{lemma}[极大无关组的判定定理]\label{lem:极大无关组的判定定理}
        设$\mathrm{rank}\left(\bm{A}\right)=r$,
    $\bm{A}=\left\{\bm{\beta}_1,\bm{\beta}_2,\cdots,
    \bm{\beta}_n\right\}$,当向量组$\left\{
    \bm{\beta}_{i_1},\bm{\beta}_{i_2},\cdots,
    \bm{\beta}_{i_r}\right\}$满足下面之一：
    \begin{enumerate}[label=\textup{(\arabic*)}]
    \item $\bm{\beta}_{i_1},\bm{\beta}_{i_2},\cdots,
    \bm{\beta}_{i_r}$线性无关.
    \item $\forall \bm{\beta}_j$都是
    $\bm{\beta}_{i_1},\bm{\beta}_{i_2},\cdots,
    \bm{\beta}_{i_r}$的线性组合.
    \end{enumerate}
    那么$\bm{\beta}_{i_1},\bm{\beta}_{i_2},\cdots,
    \bm{\beta}_{i_r}$是$\bm{A}$的极大无关组.
\end{lemma}
\end{green}
\begin{green}
\begin{lemma}[阶梯形与秩]\label{lem:阶梯形与秩}
    设$\bm{A}$为阶梯形矩阵,
    $a_{1k_1},a_{2k_2},\cdots,
    a_{rk_r}$是阶梯点,则$\mathrm{rank}\left(\bm{A}
    \right)=r=\bm{A}$的非零行数,
    且阶梯点所在列向量构成$\bm{A}$的列向量的极大无关组.
\end{lemma}
\begin{Proof}设阶梯形矩阵
\[
\bm{A}=
\begin{bmatrix}
    0&\cdots&0&a_{1k_1}&*&*&\cdots&&&\\
    0&\cdots&\cdots&0&a_{2k_2}&*&*&\cdots&&\\
    \cdots&\cdots&\cdots&\cdots&\cdots&&&&&\\
    0&\cdots&\cdots&\cdots&0&a_{rk_r}&*&*&\cdots&\\
    &&&&&&&&&\\
\end{bmatrix}
\]
明显可以化为相抵标准型.\[
\begin{pmatrix}
        \bm{I}_r & \bm{O} \\
        \bm{O} & \bm{O}
    \end{pmatrix}
\]
那么得到秩.然后将阶梯点所在列向量取出构成新矩阵
\[
\bm{B}=
\begin{bmatrix}
    a_{1k_1}&*&*&*\\
    &a_{2k_2}&*&*\\
    &&\ddots&*\\
    &&&a_{rk_r}\\
    &&&\\
\end{bmatrix}
\]
该矩阵仍为阶梯形,$r$个列向量线性无关.由\cref{lem:极大无关组的判定定理}知此为极大无关组.

证毕. 
\end{Proof}
\end{green}
\subsection{求秩相关问题}
\begin{brown}
\begin{example}
    常见的三类问题为：

\paragraph{Case\,1}
求矩阵$\bm{A}$的秩及列向量极大无关组的方法.
\begin{enumerate}[label=\textup{(\arabic*)}]
    \item 施加初等行变换化为阶梯形矩阵$\bm{B}
    ,b_{1k_{1}},b_{2k_2},\cdots,b_{rk_{r}}$
        为阶梯点
    \item $\mathrm{r}\left(\bm{A}\right)
    =\mathrm{r}\left(\bm{B}\right)
    =\bm{B}$的非零行数$ = r$
    \item $\bm{A} = \left(\bm{\beta}_1,\bm{\beta}_2,
    \cdots,\bm{\beta}_n\right)$
        的极大无关组为
    $\left(
        \bm{\beta}_{k_1},\cdots,\bm{\beta}_{k_r}
        ,\bm{\beta}_{k_r}
    \right)$
\end{enumerate}

\paragraph{Case\,2}
行、列向量组秩的计算及线性关系判定.
\begin{enumerate}[label=\textup{(\arabic*)}]
    \item 把向量组拼成矩阵求秩$\mathrm{r}\left(\bm{A}
    \right)$
    \item $\displaystyle
    \mathrm{r}\left(\bm{A}\right)
    \begin{cases*}
    = \text{向量个数}\Longrightarrow \text{
        线性无关
    }    \\
    < \text{向量个数}\Longrightarrow 
    \text{线性相关（阶梯形矩阵至少有一行零）}
    \end{cases*}
    $
\end{enumerate}

\paragraph{Case\,3}
求行、列向量组的一组极大无关组.
\begin{enumerate}[label=\textup{(\arabic*)}]
    \item 按照列向量（行向量需转置）方式拼成矩阵
    $\bm{A}$
    \item 求出秩和极大无关组
\end{enumerate}
\end{example}
\end{brown}
\subsection{子式判别法}
\begin{green}
\begin{lemma}[截断的线性关系]\label{lem:截断的线性关系}
    设$\bm{A}\in M_{m\times n}\left(\mathbb{K}\right),
    \bm{A}=\left\{
        \bm{\beta}_1,\bm{\beta}_2,\cdots,\bm{\beta}_n
    \right\}$线性相关,且$0 < r < m$,
    定义截断$\tau_{\leqslant r}\bm{\beta}=\left(
        a_1,a_2,\cdots,a_r
    \right)
    $,那么这些截断组成的向量组也是线性相关的即
    \[\left\{\tau_{\leqslant r}\bm{\beta}_1,\tau_{\leqslant r}\bm{\beta}_2,\cdots,
    \tau_{\leqslant r}\bm{\beta}_n\right\}\]线
    性相关.
\end{lemma}
\end{green}
\begin{formal}
\begin{theorem}[子式判别法]\label{thm:子式判别法}
    设$\bm{A} \in M_{m \times n}
    \left(\mathbb{K}\right)$
    \[
        \mathrm{r}\left(\bm{A}\right)
        = r \Longleftrightarrow 
        \bm{A}\text{有一个非零}
        r\text{阶子式且所有}r+1
        \text{阶子式均为零}
    \]
\end{theorem}
\begin{Proof}
先证必要性,$\mathrm{r}\left(\bm{A}\right)=r$,不妨设
前$r$行线性无关.取出为
\[
\bm{B}=\begin{pmatrix}
    \bm{\alpha}_1\\
    \bm{\alpha}_2\\
    \vdots\\
    \bm{\alpha}_r
\end{pmatrix}
\Longrightarrow \mathrm{r}\left(\bm{B}\right)=\bm{B}\text{行秩}=r
\]
那么列秩也等于$r$.不妨设前$r$列线性无关.
取出为
\[
\bm{C}
=
\begin{pmatrix}
    a_{11}&a_{12}&\cdots&a_{1r}\\
    a_{21}&a_{22}&\cdots&a_{2r}\\
    \vdots&\vdots&&\vdots\\
    a_{r1}&a_{r2}&\cdots&a_{rr}
\end{pmatrix}\Longrightarrow \mathrm{r}\left(\bm{C}\right)=r
\]
故$\bm{C}$满秩,那么它非异.于是证明了存在这样一个非零
$r$阶子式.

再来证明$r+1$阶子式为零,不妨证
\[
\bm{A}\begin{pmatrix}
    1&2&\cdots&r&r+1\\
    1&2&\cdots&r&r+1
\end{pmatrix}=0
\]
因为$\mathrm{r}\left(\bm{A}\right)=r$,所以
$\bm{\alpha}_1,\bm{\alpha}_{2},\cdots,\bm{\alpha}_{r+1}$线
性相关.设$\bm{\alpha}=\left(a_1,a_2,\cdots,a_n\right)\in \mathbb{K}_n$,定义一个切断
$\tau _{\leqslant r}\bm{\alpha}=\left(a_1,a_2,\cdots,
a_r\right)$,那么由\cref{lem:截断的线性关系}得知
\[
\tau_{\leqslant r+1}\bm{\alpha}_1,\tau_{\leqslant r+1}\bm{\alpha}_2,
\cdots,\tau_{\leqslant r+1}\bm{\alpha}_{r+1}
\]
线性相关.
又
\[
\begin{pmatrix}
   \tau_{\leqslant r+1}\bm{\alpha}_1\\
   \tau_{\leqslant r+1}\bm{\alpha}_2\\
   \vdots\\
\tau_{\leqslant r+1}\bm{\alpha}_{r+1} 
\end{pmatrix}=
\begin{pmatrix}
    a_{11}&a_{12}&\cdots&a_{1,r+1}\\
    a_{21}&a_{22}&\cdots&a_{2,r+1}\\
    \vdots&\vdots&&\vdots\\
    a_{r+1,1}&a_{r+1,2}&\cdots&a_{r+1,r+1}
\end{pmatrix}
\]
这个矩阵缺秩.
于是
\[
\bm{A}\begin{pmatrix}
    1&2&\cdots&r&r+1\\
    1&2&\cdots&r&r+1
\end{pmatrix}=\left|\bm{C}\right|=0
\]
从而必要性证毕,下面证明充分性.

因为$r+1$阶子式全为零,用\cref{thm:Laplace定理}Laplace定理可以证明所有
高于$r$阶的子式全为零.设$\mathrm{r}\left(\bm{A}\right)=t$,
由上可知,$\bm{A}$有一个非零$t$阶子式,高于$t$阶子式均为零.

首先,$t>r$时,显然与所有高于$r$阶子式为零矛盾.$t<r$类
似.于是$t=r$,证毕.
\end{Proof}
\end{formal}
\subsection{关于秩的式子}
\begin{formal}
    \begin{definition}[零秩]\label{def:零秩}
        $\mathrm{r}\left(\bm{A}\right)= 0
\Longleftrightarrow 
\bm{A} = \bm{O}
$
    \end{definition}
\end{formal}

下面给出一些关于矩阵秩的公式、不等式.
\begin{formal}
\begin{theorem}[分块三角和分块对角阵]\label{thm:分块三角和分块对角阵}
经常用于构造.

$(1)$分块对角矩阵的秩
\[
\bm{C} = \begin{pmatrix}
    \bm{A} & \bm{O} \\
    \bm{O} & \bm{B}
\end{pmatrix}
\Longrightarrow \mathrm{r}\left(\bm{C}
\right) = \mathrm{r}\left(\bm{A}
\right)+\mathrm{r}\left(\bm{B}
\right)
\]
\begin{Proof}
设存在非异阵$\bm{P}_1$、$\bm{P}_2$、$\bm{Q}_1$、
$\bm{Q}_2$使得
\begin{align*}
    \bm{P}_1\bm{AQ}_1=\begin{pmatrix}
        \bm{I}_{r_1}&\bm{O}\\
        \bm{O}&\bm{O}
    \end{pmatrix}\\
    \bm{P}_2\bm{BQ}_2=\begin{pmatrix}
        \bm{I}_{r_2}&\bm{O}\\\bm{O}&\bm{O}
    \end{pmatrix}
\end{align*}
于是
\begin{align*}
\begin{pmatrix}
    \bm{P}_1&\bm{O}\\\bm{O}&\bm{Q}_1
\end{pmatrix}
\begin{pmatrix}
    \bm{A}&\bm{O}\\
    \bm{O}&\bm{B}
\end{pmatrix}
\begin{pmatrix}
    \bm{P}_2&\bm{O}\\\bm{O}&\bm{Q}_2
\end{pmatrix}
&=\begin{pmatrix}
    \bm{P}_1\bm{AQ}_1&\bm{O}\\
    \bm{O}&\bm{P}_2\bm{BQ}_2
\end{pmatrix}\\
&=\begin{pmatrix}
    \bm{I}_{r_1}&\bm{O}&\bm{O}&\bm{O}\\
    \bm{O}&\bm{O}&\bm{O}&\bm{O}\\
    \bm{O}&\bm{O}&\bm{I}_{r_2}&\bm{O}\\
    \bm{O}&\bm{O}&\bm{O}&\bm{O}
\end{pmatrix}
\end{align*}
由于分块初等变换不改变分块矩阵的秩,上面的矩阵化为
\[
\begin{pmatrix}
    \bm{I}_{r_1+r_2}&\bm{O}\\
    \bm{O}&\bm{O}
\end{pmatrix}
\]
于是得到$\mathrm{r}\left(\bm{C}
\right)=r_1+r_2$.得证.
\end{Proof}
$(2)$分块三角矩阵的秩
\[
\bm{C}=
\begin{pmatrix}
    \bm{A} & \bm{D} \\
    \bm{O} & \bm{B}
\end{pmatrix}
\text{或}
\begin{pmatrix}
    \bm{A} & \bm{O} \\
    \bm{D} & \bm{B}
\end{pmatrix}
\Longrightarrow 
\mathrm{r}\left(\bm{C}
\right) \geqslant  \mathrm{r}\left(\bm{A}
\right)+\mathrm{r}\left(\bm{B}
\right)
\]
等号当且仅当矩阵方程
\[
\bm{AX}+\bm{BY} = \bm{D}
\]
有解时取到
\end{theorem}
\begin{Proof}
仿照作
\begin{align*}
\begin{pmatrix}
    \bm{P}_1&\bm{O}\\\bm{O}&\bm{Q}_1
\end{pmatrix}
\begin{pmatrix}
    \bm{A}&\bm{D}\\
    \bm{O}&\bm{B}
\end{pmatrix}
\begin{pmatrix}
    \bm{P}_2&\bm{O}\\\bm{O}&\bm{Q}_2
\end{pmatrix}
&=\begin{pmatrix}
    \bm{P}_1\bm{AQ}_1&\bm{P}_1\bm{DQ}_2\\
    \bm{O}&\bm{P}_2\bm{BQ}_2
\end{pmatrix}\\
&=\begin{pmatrix}
    \bm{I}_{r_1}&\bm{O}&\bm{D}_{11}&\bm{D}_{12}\\
    \bm{O}&\bm{O}&\bm{D}_{21}&\bm{D}_{22}\\
    \bm{O}&\bm{O}&\bm{I}_{r_2}&\bm{O}\\
    \bm{O}&\bm{O}&\bm{O}&\bm{O}
\end{pmatrix}\\
&\longrightarrow
\begin{pmatrix}
    \bm{I}_{r_1}&\bm{O}&\bm{O}&\bm{O}\\
    \bm{O}&\bm{O}&\bm{O}&\bm{D}_{22}\\
    \bm{O}&\bm{O}&\bm{I}_{r_2}&\bm{O}\\
    \bm{O}&\bm{O}&\bm{O}&\bm{O}
\end{pmatrix}\\
&\longrightarrow
\begin{pmatrix}
    \bm{I}_{r_1}&\bm{O}&\bm{O}&\bm{O}\\
    \bm{O}&\bm{D}_{22}&\bm{O}&\bm{O}\\
    \bm{O}&\bm{O}&\bm{I}_{r_2}&\bm{O}\\
    \bm{O}&\bm{O}&\bm{O}&\bm{O}
\end{pmatrix}
\end{align*}
得证.取等条件为$\bm{D}_{22}=\bm{O}$,经过优化为
当且仅当矩阵方程
\[
\bm{AX}+\bm{YB}=\bm{D}
\]
有解.
\end{Proof}
\end{formal}
\newpage
\begin{formal}
\begin{theorem}[秩与初等变换]\label{thm:秩与初等变换}
矩阵乘以可逆阵（即做初等变换）
不改变秩；分块矩阵在分块初等变换
下也不改变秩
\end{theorem}
\end{formal}
\begin{formal}
\begin{theorem}[秩的降阶公式]\label{thm:秩的降阶公式}
    设
    \[
\bm{M} = \begin{pmatrix}
    \bm{A} & \bm{B} \\
    \bm{C} & \bm{D} 
\end{pmatrix}
\]
那么有
%\setcounter{equation}{0}%重置计数器,方程下计数器为equation
\begin{enumerate}[label=\textup{(\arabic*)}]
   \item $\bm{A}$可逆则
$\mathrm{r}\left(\bm{M}\right)
= \mathrm{r}\left(\bm{A}\right)
+ \mathrm{r}\left(
    \bm{D} - \bm{C}\bm{A}^{-1}
    \bm{B}
\right) $
\item $\bm{D}$可逆则$
\mathrm{r}\left(\bm{M}\right)
=\mathrm{r}\left(\bm{D}\right)
+\mathrm{r}\left(
    \bm{A}-
    \bm{B}\bm{D}^{-1}\bm{C}
\right)$
\item $\bm{A}$、$\bm{D}$均可逆则$
\mathrm{r}\left(\bm{A}\right)
+ \mathrm{r}\left(
    \bm{D} - \bm{C}\bm{A}^{-1}
    \bm{B}
\right)=
\mathrm{r}\left(\bm{D}\right)
+\mathrm{r}\left(
    \bm{A}-
    \bm{B}\bm{D}^{-1}\bm{C}
\right)$
\end{enumerate}
\end{theorem}
\begin{Proof}
仅证$(1)$.
\begin{align*}
    \bm{M}&\longrightarrow\begin{pmatrix}
        \bm{A}&\bm{B}
        \\
        \bm{O}&\bm{D}-\bm{CA}^{-1}\bm{B}
    \end{pmatrix}
    \\
    &\longrightarrow
    \begin{pmatrix}
        \bm{A}&\bm{O}\\
        \bm{O}&\bm{D}-\bm{CA}^{-1}\bm{B}
    \end{pmatrix}\qedhere
\end{align*}
\end{Proof}
\end{formal}

\begin{formal}
\begin{theorem}[运算的秩]\label{thm:运算的秩}
包括加、减、数乘、分块、转置
\begin{enumerate}[label=\textup{(\arabic*)}]
    \item  $\mathrm{r}\left(k\bm{A}\right)
    =\mathrm{r}\left(\bm{A}\right)
    \left(k \neq 0\right)$
    \item $\mathrm{r}\left(\bm{A}+\bm{B}
    \right)\leqslant \mathrm{r}
    \left(\bm{A}\right) + \mathrm{r}
    \left(\bm{B}\right)$
    \item $\left|
        \mathrm{r}\left(\bm{A}\right)
        -
        \mathrm{r}\left(\bm{B}\right)
    \right|
    \leqslant 
    \mathrm{r}\left(\bm{A}
    -\bm{B}\right)
    \leqslant 
    \mathrm{r}\left(\bm{A}\right)+
    \mathrm{r}\left(\bm{B}\right)$
    \item $\mathrm{r}
    \left (
        \begin{array}{c:c}
            \begin{matrix}
                \bm{A}
            \end{matrix}
            &
            \begin{matrix}
                \bm{B}
            \end{matrix}
        \end{array}
    \right )\leqslant
    \mathrm{r}\left(\bm{A}\right)+
    \mathrm{r}\left(\bm{B}\right)$
    \item $\mathrm{r}
    \left ( 
    \begin{array}{c}
        \bm{A} \\
        \hdashline
        \bm{B}
    \end{array}
    \right )
    \leqslant
    \mathrm{r}\left(\bm{A}\right)+
    \mathrm{r}\left(\bm{B}\right)$
    \newpage
    \item $
    \bm{A} \in M_{m \times n}\left(
        \mathbb{K}
    \right) ,\bm{B}\in 
    M_{n \times p}
    \left(\mathbb{K}\right)$,则\[
    \mathrm{r}\left(\bm{A}\right)+
    \mathrm{r}\left(\bm{B}\right)
    - n\leqslant 
    \mathrm{r}\left(
        \bm{AB}\right)
        \leqslant 
    \min\left\{
    \mathrm{r}\left(\bm{A}\right),
    \mathrm{r}\left(\bm{B}\right)
    \right\}\]
    左侧
    称为\textup{Sylvester}不等式
    \item $\mathrm{r}\left(\bm{A}\right)
    +\mathrm{r}\left(\bm{I}_n
    +\bm{A}\right)\geqslant 
    n$
    \item 对于实矩阵$
    \mathrm{r}\left(\bm{AA}^*\right)
    =\mathrm{r}\left(\bm{A}^*\bm{A}
    \right)=\mathrm{r}\left(\bm{A}\right)
    $
\end{enumerate}
\end{theorem}
\begin{Proof}
$(6)$先证明上界,先证$\mathrm{r}\left(\bm{AB}\right)\leqslant \mathrm{r}\left(\bm{B}
\right)$.

令$\bm{B}=\left(
    \bm{\beta}_1,\bm{\beta}_2,\cdots,\bm{\beta}_p
\right)$
,其极大无关组为$\left\{
    \bm{\beta}_{i_1},\bm{\beta}_{i_2},\cdots,\bm{\beta}_{i_r}
\right\}$,可以断言$\forall 1\leqslant j\leqslant p$,
$\bm{A\beta}_j$是$
    \bm{A\beta}_{i_1},\bm{A\beta}_{i_2},
    \cdots,\bm{A\beta}_{i_r}
$的线性组合.前者正是$\bm{AB}$的列向量,于是
$\mathrm{r}\left(\bm{AB}\right)\leqslant r=\mathrm{r}
\left(\bm{B}\right).$一方面证毕.

另一方面,考虑转置,$\mathrm{r}\left(\bm{AB}\right)=
\mathrm{r}\left(\bm{B}^{\prime}\bm{A}^{\prime}\right)\leqslant 
\mathrm{r}\left(\bm{A}^{\prime}\right)=\mathrm{r}\left(\bm{A}
\right)$.

综上,上界证毕.下面证明下界Sylvester不等式.

作
\begin{align*}
    \begin{pmatrix}
        \bm{A}&\bm{O}\\
        -\bm{I}_n&\bm{B}
    \end{pmatrix}
    &\longrightarrow
    \begin{pmatrix}
        \bm{O}&\bm{AB}\\
        -\bm{I}_n&\bm{B}
    \end{pmatrix}\\
    &\longrightarrow
    \begin{pmatrix}
        \bm{O}&\bm{AB}\\
        -\bm{I}_n&\bm{O}
    \end{pmatrix}\\
    &\longrightarrow
    \begin{pmatrix}
        \bm{AB}&\bm{O}\\
        \bm{O}&\bm{I}_n
    \end{pmatrix}
\end{align*}
于是
\[
\mathrm{r}\left(\bm{A}\right)+\mathrm{r}\left(\bm{B}\right)\leqslant
\mathrm{r}\left(\bm{AB}\right)+n\qedhere
\]
\end{Proof}
\end{formal}
\begin{formal}
\begin{theorem}[伴随阵的秩]\label{thm:伴随阵的秩}
    对于矩阵$\bm{A}\in M_n\left(\mathbb{K}\right)$,其伴随矩阵$\bm{A}^*$的秩为\[
    \mathrm{r}\,\left(\bm{A}^*\right)=\begin{cases*}
    n&,$\mathrm{r}\,\left(\bm{A}\right)=n$\\
    1&,$\mathrm{r}\,\left(\bm{A}\right)=n-1$\\
    0&,$\mathrm{r}\,\left(\bm{A}\right)<n-1$
    \end{cases*}
    \]
\end{theorem}
\begin{Proof}
    根据\cref{prop:伴随阵性质总结}即得.
\end{Proof}
\end{formal}
\begin{brown}
\begin{example}下面给出几种特殊矩阵的秩
    \paragraph{幂等阵}幂等阵$\bm{A}$
满足$\bm{A}^2=\bm{A}$
,其秩满足
\[
    \mathrm{r}\left(\bm{A}\right)
    +\mathrm{r}\left(\bm{I}_n
    -\bm{A}\right) = 
    n
\]
\begin{Proof}
考虑
\begin{align*}
    \begin{pmatrix}
        \bm{A}&\bm{O}\\
        \bm{O}&\bm{I}_n-\bm{A}
    \end{pmatrix}&\longrightarrow
    \begin{pmatrix}
        \bm{A}&\bm{O}\\
        \bm{A}&\bm{I}_n-\bm{A}
    \end{pmatrix}\\
    &\longrightarrow
    \begin{pmatrix}
        \bm{A}&\bm{A}\\
        \bm{A}&\bm{I}_n
    \end{pmatrix}\\
    &\longrightarrow
    \begin{pmatrix}
        \bm{A}-\bm{A}^2&\bm{O}\\
        \bm{A}&\bm{I}_n
    \end{pmatrix}\\
    &\longrightarrow
    \begin{pmatrix}
        \bm{A}-\bm{A}^2&\bm{O}\\
        \bm{O}&\bm{I}_n
    \end{pmatrix}\\
    &\longrightarrow
    \begin{pmatrix}
        \bm{I}_n&\bm{O}\\\bm{O}&\bm{O}
    \end{pmatrix}
\end{align*}
证毕.
\end{Proof}
\paragraph{对合阵}对合阵$\bm{A}$满足
$\bm{A}^2 = \bm{I}_n$,其秩满足
\[
    \mathrm{r}\left(
        \bm{I}_n - \bm{A}\right)
    +\mathrm{r}\left(\bm{I}_n
    -\bm{A}\right) = 
    n
\]
\end{example}
\end{brown}

\newpage
\section{坐标向量}
\subsection{映射}
\begin{green}
\begin{lemma}[线性表出唯一性]\label{lem:线性表出唯一性}
    设$\left\{
        \bm{e}_1,\bm{e}_2,\cdots,
        \bm{e}_n
    \right\}$为$V_{
        \mathbb{K}
    }$的一组基且对于
    $\bm{\alpha}\in V$有
    \[
    \bm{\alpha} =
    a_1\bm{e}_1+a_2\bm{e}_2+\cdots+a_n\bm{e}_n
    =
    b_1\bm{e}_1 + b_2\bm{e}_2+\cdots + 
    b_n\bm{e}_n
    \]
    则\[
    a_i = b_i \quad
    \left(i = 1,2,\cdots,n\right)
    \]
\end{lemma}
\begin{Proof}
移项,由线性无关即得.
\end{Proof}
\end{green}
\begin{formal}
\begin{definition}[坐标向量]\label{def:坐标向量}
    固定$V_{\mathbb{K}}$的一组基
    $\left\{\bm{e}_1,\bm{e}_2,\cdots
    ,\bm{e}_n\right\}$及其顺序,定义
    映射$\varphi $
    \begin{align*}
        \varphi :
        V_{\mathbb{K}} & 
        \longrightarrow 
        \mathbb{K}^n \\
        \bm{\alpha} = \sum_{i = 1}
        ^{n}a_i\bm{e}_i 
        & \longmapsto \begin{pmatrix}
            a_1 \\
            a_2 \\
            \vdots \\
            a_n
        \end{pmatrix}
    \end{align*}
其中$\left(a_1,a_2,\cdots,a_n\right)^{\prime}$
是$\bm{\alpha}$在$\left\{
    \bm{e}_1,\bm{e}_2,\cdots,\bm{e}_n
\right\}$下的坐标向量.
容易证明$\varphi$是双射.
\end{definition}
\end{formal}
\subsection{同构}
\begin{formal}
\begin{definition}[线性同构]\label{def:线性同构}
    设$V_{\mathbb{K}},U_{\mathbb{K}}$
    ,若存在一个双射$\varphi
    : V \longrightarrow U$使得
    $\forall \bm{\alpha},\bm{\beta}
     \in V,k \in \mathbb{K}$,满足
     \[
     \varphi \left(\bm{\alpha}
     +\bm{\beta}
     \right) =
     \varphi \left(\bm{\alpha}
     \right)+\varphi \left(
        \bm{\beta}
     \right);
     \varphi \left(k\bm{\alpha}
     \right)=k\varphi \left(
        \bm{\alpha}
     \right)
     \]
     即$\varphi $保持了加法和
     数乘也就是保持了线性组合,
     则称$\bm{\varphi}:V\longrightarrow U$
     为（线性）同构,简称$V$同构于$U$,
     记作$V\cong U$.
\end{definition}
\end{formal}
\begin{formal}
\begin{theorem}[线性同构的关系]\label{thm:线性同构的关系}
    同构关系是一种等价关系.
\end{theorem}
\end{formal}
\begin{formal}
\begin{theorem}[同构判定]\label{thm:同构判定}
    $V_{\mathbb{K}},U_{\mathbb{K}}$
    同构
等价于
    \[
    \dim_{\mathbb{K}}\,V
    = \dim_{\mathbb{K}}\,U
    \]
\end{theorem}
\end{formal}
\begin{formal}
\begin{theorem}[与行、列向量空间同构]\label{thm:与行、列向量空间同构}
    $\mathbb{K}$上
    任一$n$维线性空间$V$
    均与$n$维列向量空间$\mathbb{K}^n$同构.
\end{theorem}
\begin{Proof}
设$\bm{\alpha}$、$\bm{\beta}\in V$,
$\displaystyle
\bm{\alpha}=\sum_{i=1}^{n}a_i\bm{e}_i$,
$\displaystyle
\bm{\beta}=\sum_{i=1}^{n}b_i\bm{e}_i$,则
\[
\bm{\varphi}\left(\bm{\alpha}\right)=\begin{pmatrix}
    a_1\\a_2\\\vdots\\a_n
\end{pmatrix},\bm{\varphi}\left(\bm{\beta}\right)=
\begin{pmatrix}
    b_1\\b_2\\\vdots\\b_n
\end{pmatrix}
\]
又$\displaystyle
\bm{\alpha}+\bm{\beta}=\sum_{i=1}^{n}\left(a_i+b_i\right)\bm{e}_i$,于是
\[
\bm{\varphi}\left(\bm{\alpha}+\bm{\beta}\right)=
\begin{pmatrix}
    a_1+b_1\\a_2+b_2\\\vdots\\a_n+b_n
\end{pmatrix}=\bm{\varphi}\left(\bm{\alpha}\right)+\bm{\varphi}
\left(\bm{\beta}\right)
\]
又因为$\displaystyle
k\bm{\alpha}=\sum_{i=1}^{n}ka_i\bm{e}_i$,于是
\[
\bm{\varphi}\left(k\bm{\alpha}\right)=\begin{pmatrix}
    ka_1\\ka_2\\\vdots\\ka_n
\end{pmatrix}=k\bm{\varphi}\left(\bm{\alpha}\right)
\qedhere
\]
\end{Proof}
\end{formal}
\begin{formal}
\begin{theorem}[线性同构之间线性结构的相同]\label{thm:线性同构之间线性结构的相同}
    线性同构保持向量组的线性关系.

    \begin{enumerate}[label=\textup{(\arabic*)}]
        \item $\varphi :V \longrightarrow 
        U$为线性空间的同构,则$
        \varphi \left(\bm{0}\right)
        = \bm{0}$
        \item $\varphi$ 将向量的线性
        组合映射为对应向量的线性组合,将
        线性相关
        的向量映射为线性相关的向量,将线性无关
        的向量映射为线性无关的向量
    \end{enumerate}
\end{theorem}
\begin{Proof}

$(1)$首先
\[
\bm{\varphi}\left(\bm{0}+\bm{0}\right)=\bm{\varphi}\left(\bm{0}\right)+\bm{\varphi}\left(
    \bm{0}\right)=\bm{\varphi}\left(\bm{0}\right)
\]
进一步
\[
\bm{\varphi}^{-1}\left(\bm{0}\right)=
\left\{\bm{\alpha}\in V\big| \bm{\varphi}\left(\bm{\alpha}\right)=\bm{0}\right\}
=\left\{\bm{0}\right\}
\]

$(2)$先证线性组合,设$\bm{\beta}=\lambda_1\bm{\alpha}_1+\lambda_2\bm{\alpha}_2+\lambda_m\bm{\alpha}_m$,则
\begin{align*}
    \bm{\varphi}\left(\bm{\beta}\right)&=\bm{\varphi}\left(\lambda_1\bm{\alpha}_1+\lambda_2\bm{\alpha}_2+\cdots+\lambda_m\bm{\alpha}_m\right)\\
    &=\lambda_1\bm{\varphi}\left(\bm{\alpha}_1\right)+\lambda_2\bm{\varphi}\left(\bm{\alpha}_2\right)+\cdots+\lambda_m\bm{\varphi}\left(
        \bm{\alpha}_m\right)
\end{align*}

然后证明线性相关,设$\lambda_1\bm{\alpha}_1+\lambda_2\bm{\alpha}_2+\cdots+\lambda_m\bm{\alpha}_m=\bm{0}$,$\lambda_i$不全为零,
于是
\begin{align*}
    \bm{0}=\bm{\varphi}\left(\bm{0}\right)&=\bm{\varphi}\left(\lambda_1\bm{\alpha}_1+\lambda_2\bm{\alpha}_2+\cdots+\lambda_m\bm{\alpha}_m\right)\\
    &=\lambda_1\bm{\varphi}\left(\bm{\alpha}_1\right)+\lambda_2\bm{\varphi}\left(\bm{\alpha}_2\right)+\cdots+\lambda_m\bm{\varphi}\left(\bm{\alpha}_m
    \right)
\end{align*}

最后是线性无关,设$\bm{\alpha}_1,\bm{\alpha}_2,\cdots,\bm{\alpha}_m$线性无关,则设
$\lambda_1\bm{\varphi}\left(\bm{\alpha}_1\right)+\lambda_2\bm{\varphi}\left(\bm{\alpha}_2\right)+\cdots+\lambda_m\bm{\varphi}\left(\bm{\alpha}_m
    \right)=\bm{0}$即
    \[
    \bm{0}=\bm{\varphi}\left(\lambda_1\bm{\alpha}_1+\lambda_2\bm{\alpha}_2+\cdots
    +\lambda_m\bm{\alpha}_m\right)
    \]
    由\[
\bm{\varphi}^{-1}\left(\bm{0}\right)=
\left\{\bm{\alpha}\in V\big| \bm{\varphi}\left(\bm{\alpha}\right)=\bm{0}\right\}
=\left\{\bm{0}\right\}
\]
得到
\[
\lambda_1\bm{\alpha}_1+\lambda_2\bm{\alpha}_2+\cdots
    +\lambda_m\bm{\alpha}_m=\bm{0}
\]
又因为$\bm{\alpha}_1,\bm{\alpha}_2,\cdots,\bm{\alpha}_m$线性无关,故
$\lambda_1=\lambda_2=\cdots=\lambda_m=0.$

证毕.
\end{Proof}
\end{formal}
\begin{formal}
\begin{theorem}[坐标向量的同构]\label{thm:坐标向量的同构}
    设$\varphi :
    V \longrightarrow 
    \mathbb{K}^n$,$\bm{\alpha}_1,\bm{\alpha}_2
    ,\cdots,\bm{\alpha}_m\in V$,
    作
    \[
    \widetilde{\bm{\alpha}}_i:=
    \varphi\left(\bm{\alpha}_i
    \right) 
    \]
    则$\left\{\bm{\alpha}_1,\bm{\alpha}_2
    ,\cdots,\bm{\alpha}_m
    \right\}$
    与$\left\{\widetilde{\bm{\alpha}
    }_1,\widetilde{\bm{\alpha}}_2,\cdots,\widetilde{
        \bm{\alpha}
    }_m  \right\}$的秩相同.

    甚至可以断言,若$\left\{
        \bm{\alpha}_{i_1},
        \bm{\alpha}_{i_2},
        \cdots
        ,\bm{\alpha}_{i_r}
    \right\}$为
    $\left\{\bm{\alpha}_1,\bm{\alpha}_2
    ,\cdots,\bm{\alpha}_m
    \right\}$的极大线性无关组,那么
    $\left\{\widetilde{\bm
    {\alpha}
    }_{i_1},
    \widetilde{\bm{\alpha}}_{i_2},
    \cdots,\widetilde{
        \bm{\alpha}
    }_{i_r}  \right\}$
    是$\left\{\widetilde{\bm{\alpha}
    }_1,
    \widetilde{\bm{\alpha}}_2,
    \cdots,\widetilde{
        \bm{\alpha}
    }_m \right\}$的极大无关组.

    即秩相同且极大无关组指标相同.
\end{theorem}
\begin{Proof}
根据\cref{thm:与行、列向量空间同构}立得.
\end{Proof}
\end{formal}
\begin{red}
\begin{remark}
    以上定理\cref{thm:坐标向量的同构}是
    一个非常重要的结论,它说明了
    同构保持了向量组的线性关系,这
    也是线性代数中一个非常重要的
    性质,在后面的许多定理证明中
    都会用到这个性质.然而更多时候你并未意识到这一点.
\end{remark}
\end{red}
\newpage
\section{基变换与过渡矩阵}
\subsection{定义}
\begin{formal}
\begin{definition}[过渡矩阵]\label{def:过渡矩阵}
    设$V_{\mathbb{K}}$的两组基
    $\left\{\bm{e}_1,\bm{e}_2,\cdots,
    \bm{e}_n\right\}$,$
    \left\{\bm{f}_1,\bm{f}_2
    ,\cdots,\bm{f}_n\right\}$
    则有\[
    \begin{cases*}
        \bm{f}_1 = a_{11}\bm{e}_1+a_{21}\bm{e}_2
        +\cdots+a_{n1}\bm{e}_n\\
        \bm{f}_2=a_{12}\bm{e}_1+a_{22}\bm{e}_2+\cdots+
        a_{n2}\bm{e}_n\\
        \qquad\cdots\cdots
        \cdots\cdots\cdots\cdots\\
        \bm{f}_n=a_{1n}\bm{e}_1+a_{2n}\bm{e}_2+\cdots
        +a_{nn}\bm{e}_n   
    \end{cases*}
    \]
    这里的系数构成一个$n$阶方阵的转置
    \[
    \bm{A} = \begin{pmatrix}
        a_{11} & a_{21}&\cdots & a_{n1} \\
        a_{12}&a_{22}&\cdots&a_{n2}\\
        \vdots & \vdots& & \vdots \\
        a_{1n} & a_{2n}&\cdots & a_{nn}
    \end{pmatrix}
    \]
    称为从基$\left\{\bm{e}_1,\bm{e}_2,\cdots,
    \bm{e}_n\right\}$到
    基$
    \left\{\bm{f}_1,\bm{f}_2,
    ,\cdots,\bm{f}_n\right\}$的过渡矩阵.
\end{definition}
\end{formal}
\begin{formal}
\begin{definition}[形式行向量]\label{def:形式行向量}
    有别于一般向量,称
    \[
    \left(\bm{\alpha}_1,\bm{\alpha}_2,\cdots
    ,\bm{\alpha}_n\right) \quad
    \left(\bm{\alpha}_i \in V_{
        \mathbb{K}
    }\right)
    \]
    是一个形式行向量.
\end{definition}
\end{formal}
\begin{formal}
\begin{definition}[形式向量的运算定义]\label{prop:形式向量的运算定义}
    类似于\cref{def:运算},定义形式向量的运算
    \begin{enumerate}[label=\textup{(\arabic*)}]
        \item 相等:
    \[
    \left(\bm{\alpha}_1,\bm{\alpha}_2,\cdots
    ,\bm{\alpha}_n\right)
    =\left(\bm{\beta}_1,\bm{\beta}_2,\cdots
    ,\bm{\beta}_n\right)
    \Longleftrightarrow 
    \bm{\alpha}_i = 
    \bm{\beta}_i,\forall
    1\leqslant i\leqslant n
    \]
    \item 加法\[
    \left(\bm{\alpha}_1,\bm{\alpha}_2,\cdots
    ,\bm{\alpha}_n\right)
    +\left(\bm{\beta}_1,\bm{\beta}_2,\cdots
    ,\bm{\beta}_n\right)
    =\left(
        \bm{\alpha}_1+\bm{\beta}_1
        ,\bm{\alpha}_1+\bm{\beta}_2,
        \cdots,\bm{\alpha}_n+
        \bm{\beta}_n
    \right)
    \]
    \item 数乘（$k\in\mathbb{K}$）\[
        k\left(\bm{\alpha}_1,\bm{\alpha}_2,\cdots
    ,\bm{\alpha}_n\right)=
    \left(k\bm{\alpha}_1,k\bm{\alpha}_2,\cdots
    ,k\bm{\alpha}_n\right)
    \]
    \item 与矩阵的乘法
    \[
    \left(\bm{\alpha}_1,\bm{\alpha}_2,\cdots
    ,\bm{\alpha}_n\right)\bm{A}_{
        n \times m
    } = 
    \left(
        \sum_{i = 1}^{n}a_{i1}\bm{\alpha}_i,
        \sum_{i=1}^{n}a_{i2}\bm{\alpha}_i
        \cdots,
        \sum_{i = 1}^{n}a_{im}\bm{\alpha}_i
    \right)
    \]
    \end{enumerate}
\end{definition}
\end{formal}
\begin{green}
\begin{corollary}[过渡矩阵的形式表出]\label{cor:过渡矩阵的形式表出}
    有了\cref{def:形式行向量}中定义的形式行向量,\cref{def:过渡矩阵}的基变换也可以写为
    \[
    \left(\bm{f}_1,\bm{f}_2,\cdots,\bm{f}_n\right)
    =\left(\bm{e}_1,\bm{e}_2,\cdots,\bm{e}_n\right)\bm{A}
    \]
其中$\bm{A}$即为过渡矩阵.
\end{corollary}
\end{green}
\begin{green}
\begin{lemma}[过渡矩阵唯一性]\label{lem:过渡矩阵唯一性}
    设$V_{\mathbb{K}}$的一组基
    $\left\{
        \bm{e}_1,\bm{e}_2,\cdots
        ,\bm{e}_n
    \right\}$,
    $\bm{A}=\left(a_{ij}\right)_{
        n \times m
    }$,$\bm{B} = 
    \left(b_{ij}\right)_{
        n \times m
    }$使得
    \[
    \left(
        \bm{e}_1,\bm{e}_2,\cdots
        ,\bm{e}_n
    \right)\bm{A}=
    \left(
        \bm{e}_1,\bm{e}_2,\cdots
        ,\bm{e}_n
    \right)\bm{B}
    \]
    则$\bm{A}=\bm{B}$.
\end{lemma}
\end{green}
\begin{red}
\begin{remark}
    给定的两组不同基下的过渡矩阵是唯一确定的.
\end{remark}
\end{red}
\begin{green}
\begin{lemma}[过渡矩阵的确定]\label{lem:过渡矩阵的确定}
    设数域$\mathbb{K}$上的$n$维线性空间$V$,取定$V$的两组基
    为$\left\{\bm{e}_1,\bm{e}_2,\cdots,
    \bm{e}_n\right\}$,且$
    \left\{\bm{f}_1,\bm{f}_2
    ,\cdots,\bm{f}_n\right\}$,$\forall
    \bm{\alpha} \in V$
    有
    \[
    \bm{\alpha}=
    \lambda _1\bm{e}_1+\lambda_2\bm{e}_2+\cdots
    +\lambda _n\bm{e}_n
    =\mu _1\bm{f}_1+\mu_2\bm{f}_2+\cdots
    +\mu _n\bm{f}_n
    \]
    则根据\cref{cor:过渡矩阵的形式表出}有
    \[
    \begin{pmatrix}
        \lambda _1 \\
        \lambda_2\\
        \vdots \\
        \lambda _n
    \end{pmatrix}
    = \begin{pmatrix}
        a_{11} & a_{12}&\cdots & a_{n1} \\
        a_{21}&a_{22}&\cdots&a_{2n}\\
        \vdots & \vdots& & \vdots \\
        a_{1n} & a_{n2}&\cdots & a_{nn}
    \end{pmatrix}
    \begin{pmatrix}
        \mu _1 \\
        \mu_2\\
        \vdots \\
        \mu_n
    \end{pmatrix}
    \]
    这是矩阵\[
    \bm{A} = \begin{pmatrix}
        a_{11} & a_{21}&\cdots & a_{n1} \\
        a_{12}&a_{22}&\cdots&a_{n2}\\
        \vdots & \vdots& & \vdots \\
        a_{1n} & a_{2n}&\cdots & a_{nn}
    \end{pmatrix}
    \]是从$\left\{\bm{e}_1,\bm{e}_2,\cdots,
    \bm{e}_n\right\}$到
    $
    \left\{\bm{f}_1,\bm{f}_2
    ,\cdots,\bm{f}_n\right\}$的过
    渡矩阵
    的充分必要条件.
\end{lemma}
\begin{Proof}
    一方面,
设
\[
    \left(\bm{f}_1,\bm{f}_2,\cdots,\bm{f}_n\right)
    =\left(\bm{e}_1,\bm{e}_2,\cdots,\bm{e}_n\right)\bm{A}
\]
因为
\begin{align*}
    \bm{\alpha}&=\left(\bm{e}_1,\bm{e}_2,\cdots,\bm{e}_n\right)\begin{pmatrix}
        \lambda_1\\\lambda_2\\\vdots\\\lambda_n
    \end{pmatrix}\\
    &=\left(\bm{f}_1,\bm{f}_2,\cdots,\bm{f}_n\right)
    \begin{pmatrix}
        \mu_1\\\mu_2\\\vdots\\\mu_n
    \end{pmatrix}\\
    &=\left(\bm{e}_1,\bm{e}_2,\cdots,\bm{e}_n\right)\bm{A}
    \begin{pmatrix}
        \mu_1\\\mu_2\\\vdots\\\mu_n
    \end{pmatrix}
\end{align*}

另一方面,考虑$\bm{f}_i$的新坐标
向量为$\left(0,0,\cdots,1,\cdots,0\right)^{\prime}\in \mathbb{K}^n$,
那么由上式得其旧坐标向量
为$
\bm{A}\left(
    0,0,\cdots,1,\cdots,0
\right)^{\prime}
$
即$\bm{A}$的第$i$列$\left(
    a_{1i},a_{2i},\cdots,a_{ni}
\right)^{\prime}$,符合过渡矩阵定义.

证毕.
\end{Proof}
\end{green}
\subsection{传递}
\begin{formal}
\begin{theorem}[过渡矩阵的可逆性与复合性质]\label{thm:过渡矩阵的可逆性与复合性质}
    设数域$\mathbb{K}$上的线性空间$V$的三
    组基为
    $\left\{
        \bm{e}_1,\bm{e}_2,\cdots
        ,\bm{e}_n
    \right\}$、$
    \left\{
        \bm{f}_1,\bm{f}_2,\cdots,
        \bm{f}_n
    \right\}$、$
    \left\{
        \bm{g}_1,\bm{g}_2,\cdots
        ,\bm{g}_n
    \right\}$,且从基$\left\{
        \bm{e}_1,\bm{e}_2,\cdots
        ,\bm{e}_n
    \right\}$到基$
    \left\{
        \bm{f}_1,\bm{f}_2,\cdots,
        \bm{f}_n
    \right\}$的过渡矩阵为
    $\bm{A}$,
    从基$
    \left\{
        \bm{f}_1,\bm{f}_2,\cdots,
        \bm{f}_n
    \right\}$到基$
    \left\{
        \bm{g}_1,\bm{g}_2,\cdots
        ,\bm{g}_n
    \right\}$的过渡矩阵为$
    \bm{B}$,
    从基$
    \left\{
        \bm{e}_1,\bm{e}_2,\cdots,
        \bm{e}_n
    \right\}$到基$
    \left\{
        \bm{g}_1,\bm{g}_2,\cdots
        ,\bm{g}_n
    \right\}$的过渡矩阵为$
    \bm{C}$
    ,则
    \begin{enumerate}[label=\textup{(\arabic*)}]
        \item 任何过渡矩阵都是可逆的
        \item $\bm{C} = \bm{AB}$
    \end{enumerate}
\end{theorem}
\newpage
\begin{Proof}$(1)$
设从基$\left\{
        \bm{f}_1,\bm{f}_2,\cdots,
        \bm{f}_n
    \right\}$到基
    $\left\{
        \bm{e}_1,\bm{e}_2,\cdots,
        \bm{e}_n
    \right\}$的过渡矩阵为
    $\bm{P}$,则
    \begin{align*}
        \left(\bm{f}_1,\bm{f}_2,\cdots,\bm{f}_n\right)=\left(\bm{e}_1,\bm{e}_2,\cdots,\bm{e}_n\right)\bm{A}\\
        \left(\bm{e}_1,\bm{e}_2,\cdots,\bm{e}_n
        \right)=\left(\bm{f}_1,\bm{f}_2,\cdots,\bm{f}_n\right)\bm{P}
    \end{align*}
    于是
    \[
    \left(\bm{e}_1,\bm{e}_2,\cdots,\bm{e}_n
        \right)=\left(\bm{e}_1,\bm{e}_2,\cdots,\bm{e}_n
        \right)\bm{AP}
    \]
    即$\bm{AP}=\bm{I}_n$,证毕.

    $(2)$显然
    \begin{align*}
        \left(\bm{f}_1,\bm{f}_2,\cdots,\bm{f}_n\right)=\left(\bm{e}_1,\bm{e}_2,\cdots,\bm{e}_n\right)\bm{A}\\
        \left(\bm{g}_1,\bm{g}_2,\cdots,\bm{g}_n\right)=\left(\bm{f}_1,\bm{f}_2,\cdots,\bm{f}_n\right)\bm{B}\\
        \left(\bm{g}_1,\bm{g}_2,\cdots,\bm{g}_n\right)=\left(\bm{e}_1,\bm{e}_2,\cdots,\bm{e}_n\right)\bm{C}
    \end{align*}
    转换立刻得到.

    于是证毕.
\end{Proof}
\end{formal}
\begin{green}
\begin{corollary}[过渡矩阵的逆阵]\label{cor:过渡矩阵的逆阵}
    设同一线性空间
    中,从基$\left\{
        \bm{e}_1,\bm{e}_2,\cdots
        ,\bm{e}_n
    \right\}$到基$
    \left\{
        \bm{f}_1,\bm{f}_2,\cdots,
        \bm{f}_n
    \right\}$的过渡矩阵
    为
    $\bm{A}$
    ,从基
    $\left\{
        \bm{f}_1,\bm{f}_2,\cdots
        ,\bm{f}_n
    \right\}$到基$
    \left\{
        \bm{e}_1,\bm{e}_2,\cdots,
        \bm{e}_n
    \right\}$的过渡矩阵
    为$\bm{B}$,
    则二者是逆阵关系
    \[
    \bm{AB} =  \bm{I}_n
    \]
\end{corollary}
\begin{Proof}
    由\cref{thm:过渡矩阵的可逆性与复合性质}立得.
\end{Proof}
\end{green}
\begin{red}
\begin{remark}
    传递或者说复合这一性质的
    重要意义在于,对于非标准基之间的过渡矩阵,
    不再需要去解出一个包含$n^2$个未定元的方程,
    而是通过与标准基之间的过渡矩阵和简单的复合表示
    ,转化为
    解矩阵方程的问题,进一步地,转换为求逆
    阵这一简单问题.
\end{remark}
\end{red}
\begin{brown}
\begin{example}
    设$\mathbb{K}$上的线性空间$V$的两组基为
    \[
    \left\{
        \bm{f}_1=\left(1,0,-1\right),\bm{f}_2=\left(2,1,1\right),\bm{f}_3=\left(1,1,1\right)
    \right\}
    \]\[\left\{
    \bm{g}_1=\left(0,1,1\right),\bm{g}_2=\left(-1,1,0\right),\bm{g}_3=\left(1,2,1\right)\right\}
    \]求基$\bm{f}$到基$\bm{g}$的过渡矩阵.
\end{example}
\begin{solution}
    设标准基为$\left\{\bm{e}_1=\left(1,0,0\right),\bm{e}_2=\left(0,1,0\right),\bm{e}_3=\left(0,0,1\right)\right\}$,则基$\bm{e}$到基$\bm{f},\bm{g}$的过渡矩阵分别为\[
    \begin{pmatrix}
        1&2&1\\
        0&1&1\\
        -1&1&1
    \end{pmatrix},\begin{pmatrix}
      0&-1&1\\
      1&1&2\\
      1&0&1  
    \end{pmatrix}
    \]并设基$\bm{f}$到基$\bm{g}$的过渡矩阵为$\bm{X}$.根据\cref{thm:过渡矩阵的可逆性与复合性质}有\[
    \bm{AX}=\bm{B}\Longrightarrow \bm{X}=\bm{A}^{-1}\bm{B}
    \]显然\[
    \bm{X}=\begin{pmatrix}
        0&1&1\\-1&-3&-2\\2&4&4
    \end{pmatrix}
    \]
\end{solution}
\end{brown}
\subsection{列形式}
行列向量的表示选择并无大碍.事实上,绝大多数书本可能更倾向于
使用形式列向量.

如果用行向量$\left(\lambda_1,\lambda_2,\cdots,\lambda_n
\right)$表示坐标向量,
定义形式列向量
\[
\begin{pmatrix}
    \bm{e}_1 \\
    \bm{e}_2\\
    \cdots \\
    \bm{e}_n
\end{pmatrix}
\]
则对于过渡矩阵有
\begin{align*}
    \begin{pmatrix}
        \bm{f}_1\\\bm{f}_2\\\vdots\\\bm{f}_n
    \end{pmatrix}=
    \bm{A}\begin{pmatrix}
        \bm{e}_1\\\bm{e}_2\\\vdots\\\bm{e}_n
    \end{pmatrix}
\end{align*}
即\[
\left(\lambda _1,\lambda_2,\cdots,
\lambda _n\right)
 = \left(\mu_1,\mu_2,\cdots
 ,\mu_n\right)\bm{A}
\]

究其根本,只是一个转置的问题.
\newpage
\section{子空间}
\subsection{子空间}
\begin{formal}
\begin{definition}[子空间]\label{def:子空间}
    设线性空间$V_{\mathbb{K}}$
    ,$V_0$为$V$的非空子集,若
    $\forall \bm{\alpha},\bm{\beta}
    \in V_0,k \in \mathbb{K}$
    \[
    \bm{\alpha}+\bm{\beta} \in V_0,k\bm{\alpha
    } \in V_0
    \]
    称$V_0$是$V$的（线性）子空间
    （准确地说是子线性空间）.
\end{definition}
\end{formal}
\begin{green}
\begin{lemma}[子线性空间]\label{lem:子线性空间}
    $V_0$在$V$的加法和数乘下是$\mathbb{K}$
    的线性空间
    \begin{align*}
        + : 
        V_0 \times V_0 
        & \longrightarrow 
        V_0 \\
        \left(
            \bm{\alpha},\bm{\beta}
        \right)
        & \longmapsto
        \bm{\alpha}+\bm{\beta}
        \\
        \cdot : \mathbb{K} \times 
        V_0 
        & \longrightarrow  
        V_0 \\
        \left(k,\bm{\alpha}
        \right) 
        & \longmapsto 
        k \cdot \bm{\alpha}
    \end{align*}
\end{lemma}
\begin{Proof}
加法交换律、结合律显然成立,又$\forall\bm{\alpha}\in V_0$,$\bm{0}=\bm{\alpha}+\left(-\bm{\alpha}\right)\in V_0$.
即零向量在$V_0$中,且就是$V$中的零向量.显然负元存在.

显然数乘的性质均满足.
\end{Proof}
\end{green}
\begin{red}
\begin{remark}
    注意子空间中的运算仍然是原空间中的运算.
\end{remark}
\end{red}
\begin{green}
\begin{corollary}[线性组合封闭]\label{cor:线性组合封闭}
    子空间对线性组合封闭即若$\bm{\alpha}_1,\bm{\alpha}_2,\cdots,\bm{\alpha}_n\in V_0$,
    $\lambda_1,\lambda_2,\cdots,\lambda_n\in \mathbb{K}$,则
    \[
    \lambda_1\bm{\alpha}_1+\lambda_2\bm{\alpha}_2+\cdots+\lambda_n\bm{\alpha}_n\in V_0
    \]
\end{corollary}
\begin{Proof}
    由\cref{lem:子线性空间}立得.
\end{Proof}
\end{green}
\begin{formal}
\begin{theorem}[平凡子空间]\label{thm:平凡子空间}
    特别地,任一线性空间$V$至少有下面
两个子空间,被称为平凡子空间.
\[
\text{平凡子空间}
\begin{cases*}
    \text{零子空间：}\left\{
        \bm{0}_V
    \right\}\\
    \text{全子空间：}
        V
\end{cases*}
\]
\end{theorem}
\end{formal}
\begin{green}
\begin{lemma}[子空间的维数]\label{lem:子空间的维数}
    对于$\mathbb{K}$上的线性空间$V$的任一子空间$V_0$
    \[
    0 \leqslant 
    \dim_{\mathbb{K}}\,V_0
    \leqslant 
    \dim_{\mathbb{K}}\,V
    \]
    特别地,非平凡子空间取到严
    格不等号.
\end{lemma}
\begin{Proof}
由\cref{thm:基扩张定理}基扩张定理,这是显然的.仅证明：$\dim_{\mathbb{K}}\,V_0=\dim_{\mathbb{K}}\,V=n$时$V=V_0$.

取$V_0$的一组基$\left\{\bm{e}_1,\bm{e}_2,\cdots,\bm{e}_n\right\}$,则
其是$V$中$n$个线性无关的向量,又$\dim_{\mathbb{K}}\,V=n$,故
这是$V$的一组基.

又$\forall \bm{\alpha}\in V$,
$\lambda_1\bm{\alpha}_1+\lambda_2\bm{\alpha}_2+\cdots+\lambda_n\bm{\alpha}_n\in V_0$,则$V_0=V$.

证毕.
\end{Proof}
\end{green}
\begin{formal}
\begin{theorem}[子空间相等]\label{thm:子空间相等}
    设$V_{\mathbb{K}}$的两个子空间$V_1,V_2$,
    则$V_1=V_2$当且仅当$\dim\,V_1=\dim\,V_2.$
\end{theorem}
\end{formal}
\subsection{交空间与和空间}
\begin{formal}
\begin{definition}[交空间与和空间]\label{def:交空间与和空间}
    设$V_1,V_2
    \subseteq V$,定义它们的交为集合的交
    \[
    V_1 \cap V_2
     = \left\{
        \bm{\alpha} \mid 
        \bm{\alpha} \in V_1,
        \bm{\alpha} \in V_2
     \right\}
    \]
    和为
    \[
    V_1 + V_2 =\left\{
        \bm{\alpha} + \bm{\beta} \mid 
        \bm{\alpha} \in V_1,
        \bm{\beta} \in V_2
    \right\}
    \]
    且二者均是$V$的子空间即交空间、和空间,
    但
    \[
    V_1\cup V_2 =\left\{
        \bm{\alpha} | 
        \bm{\alpha} \in V_1 \text{或}
        \bm{\alpha} \in V_2
    \right\}
    \]
    不是$V$的子空间（一般来说）.
\end{definition}
\begin{Proof}
先证交空间,设$\bm{\alpha},\bm{\beta}\in V_1\cap V_2$,则
$\bm{\alpha}+\bm{\beta}\in V_1$且$\bm{\alpha}+\bm{\beta}\in V_2$,故
$\bm{\alpha}+\bm{\beta}\in V_1\cap V_2$.$\forall k \in\mathbb{K}$,
$\bm{\alpha}\in V_1\cap V_2$,均有$k\bm{\alpha}\in V_1$且
$k\bm{\alpha}\in V_2$,故$k\bm{\alpha}\in V_1\cap V_2.$

再证和空间,设$\bm{\alpha},\bm{\beta}\in V_1+V_2$,且
\begin{align*}
    &\bm{\alpha}=\bm{\alpha}_1+\bm{\alpha}_2,\bm{\alpha}_1\in V_1,\bm{\alpha}_2\in V_2\\
   & \bm{\beta}=\bm{\beta}_1+\bm{\beta}_2,\bm{\beta}_1\in V_1,\bm{\beta}_2\in V_2
\end{align*}
则
\[
\bm{\alpha}+\bm{\beta}=\left(\bm{\alpha}_1+\bm{\beta}_1\right)+
\left(\bm{\alpha}_2+\bm{\beta}_2\right)\in V_1+V_2
\]
考虑$\forall k \in \mathbb{K}$,$k\bm{\alpha}=k\bm{\alpha}_1
+k\bm{\alpha}_2\in V_1+V_2$.

于是证毕.
\end{Proof}
\end{formal}
\begin{red}
\begin{remark}
    两个线性空间的交空间不可能为空,因为至少包含
    零元,也就是说,下限就是零子空间$0=\left\{\bm{0}
    \right\}$.
\end{remark}
\end{red}
\begin{green}
\begin{corollary}[交空间与和空间推广]\label{cor:交空间与和空间推广}
    设$V_1,V_2,\cdots,V_m$均为$V$的子空间,
    则它们的交空间
    \[
    V_1 \cap V_2\cap\cdots \cap V_m
    = \left\{
        \bm{\alpha} \mid 
        \bm{\alpha} \in V_1 ,
        \bm{\alpha}\in V_2,
        \cdots,
        \bm{\alpha} \in V_m
    \right\}
    \]
    与和空间
    \[
    V_1 + V_2+\cdots + V_m =
    \left\{
        \bm{\alpha}_1 + \bm{\alpha}_2+\cdots +
        \bm{\alpha}_m\mid
        \bm{\alpha}_1 \in V_1,
        \bm{\alpha}_2\in V_2,
        \cdots,
        \bm{\alpha}_m \in V_m
    \right\}
    \]
    均为$V$的子空间.
\end{corollary}
\begin{Proof}
证明同\cref{def:交空间与和空间}.
\end{Proof}
\end{green}
\subsection{张成的子空间}
\begin{formal}
\begin{definition}[张成的子空间]\label{def:张成的子空间}
    设$S$为$V_{
        \mathbb{K}
    }$的非空子集,
    记$L\left(S\right)$为
    $S$中所有可能的线性组合构成的
    集合即
    \[
    L\left(S\right)
    = 
    \left\{
        \lambda _1 \bm{\alpha}_1+\lambda_2\bm{\alpha}_2+
        +\cdots+
        \lambda _m\bm{\alpha}_m
        \Bigg |
        \begin{matrix}
            \lambda _1,\lambda_2,\cdots,
            \lambda _m \in \mathbb{K}\\
            \bm{\alpha}_1,\bm{\alpha}_2,\cdots,
            \bm{\alpha}_m \in S
        \end{matrix}
        \quad
        \left(
            m \geqslant 0\textup{但有限}
        \right)
    \right\}
    \]
    其在加法、数乘下封闭,是$V$的一个子空间,
    称为由$S$生成/张成的子空间.
\end{definition}
\end{formal}
\begin{formal}
\begin{theorem}[张成的子空间的特点]\label{thm:张成的子空间的特点}
    $S$为$V_{\mathbb{K}}$的子集,$L\left(
        S
    \right)$是$S$张成的子空间,则
    \begin{enumerate}[label=\textup{(\arabic*)}]
        \item $L\left(S\right)$是包含$S
        $的$V$的最小子空间
        \item 若$S$存在极大无关组$
        \left\{\bm{\alpha}_1,\bm{\alpha}_2,\cdots
        ,\bm{\alpha}_r\right\}$
            ,则
        \[
        L\left(S\right)
        = L\left(\bm{\alpha}_1,\bm{\alpha}_2,
        \cdots
        ,\bm{\alpha}_r\right)
        \]
        \newpage
        且$
        \left\{\bm{\alpha}_1,\bm{\alpha}_2,
        \cdots
        ,\bm{\alpha}_r\right\}$
        是$
        L\left(S
        \right)
        $的一组基即\[
        \dim_{\mathbb{K}}\,L\left(S
        \right) = \mathrm{r}\left(
        S\right)
        \]
    \end{enumerate}
\end{theorem}
\begin{Proof}
$(1)$任取子空间$V_0\supseteq S$,任取$\lambda_1\bm{\alpha}_1+\lambda_2\bm{\alpha}_2+\cdots+\lambda_m\bm{\alpha}_m\in L\left(S\right)$,
$\bm{\alpha}_i\in S$,故$\bm{\alpha}_i\in V_0$,又因为子空间
保持线性组合,于是$\lambda_1\bm{\alpha}_1+\lambda_2\bm{\alpha}_2+\cdots+\lambda_m\bm{\alpha}_m\in V_0$,
于是$L\left(S\right)\subseteq V_0.$

$(2)$因为
\[
L\left(S\right)\hookrightarrow S\hookrightarrow
\left\{\bm{\alpha}_1,\bm{\alpha}_2,\cdots,
\bm{\alpha}_r\right\}
\]
显然
\[
L\left(S\right)
        = L\left(\bm{\alpha}_1,\bm{\alpha}_2,
        \cdots
        ,\bm{\alpha}_r\right)
        \]
且$\left\{\bm{\alpha}_1,\bm{\alpha}_2,\cdots,\bm{\alpha}_r\right\}$
是$L\left(S\right)$的一组基.于是
\[
\dim_{\mathbb{K}}L\left(S\right)=r=\mathrm{r}\left(S\right)
\qedhere
\]
\end{Proof}
\end{formal}
\begin{green}
\begin{corollary}[并集张成的空间]\label{cor:并集张成的空间}
    若$V_1,V_2,\cdots,V_m$为$V$的子空间,
    则
    \[
    L\left(
        V_1 \cup V_2\cup \cdots \cup
        V_m 
    \right)=
    V_1 +V_2+\cdots
    + V_m
    \]
\end{corollary}
\begin{Proof}
考虑$\forall \bm{\alpha}=\bm{\alpha}_1+\bm{\alpha}_2+
\cdots+\bm{\alpha}_m\in V_1+V_2+\cdots+V_m$,其中
$\bm{\alpha}_i\in V_i$,从而
\[
\bm{\alpha}=\bm{\alpha}_1+\bm{\alpha}_2+
\cdots+\bm{\alpha}_m\in L\left(
    V_1 \cup V_2\cup \cdots \cup
        V_m 
\right)
\]
即有
\[
    L\left(
        V_1 \cup V_2\cup \cdots \cup
        V_m 
    \right)\supseteq
    V_1 +V_2+\cdots
    + V_m
    \]
又因为$V_i\subseteq V_1 + V_2+ \cdots +
        V_m $,则$\bm{\alpha}_i\in V_1+V_2+\cdots+V_m$,
于是
\[
V_1\cup V_2\cup \cdots\cup V_m\subseteq V_1+V_2+\cdots+V_m
\]
由最小性即得
\[
    L\left(
        V_1 \cup V_2\cup \cdots \cup
        V_m 
    \right)\subseteq
    V_1 +V_2+\cdots
    + V_m
    \]
于是有\[
    L\left(
        V_1 \cup V_2\cup \cdots \cup
        V_m 
    \right)=
    V_1 +V_2+\cdots
    + V_m
\qedhere
    \]
    \end{Proof}
\end{green}
\subsection{维数公式}
\begin{formal}
\begin{theorem}[子空间的维数公式]\label{thm:子空间的维数公式}
    设$V_1,V_2$为$V_{\mathbb{K}}$的子空间,
    则\[
    \dim_{\mathbb{K}}\,\left(V_1 +
    V_2\right)
    =
    \dim_{\mathbb{K}}\,V_1+
    \dim_{\mathbb{K}}\,V_2 -
    \dim_{\mathbb{K}}\,\left(
        V_1\cap V_2
    \right)
    \]
    特别地,$V_1 \cap V_2 =0=\left\{
        \bm{0}
    \right\}$时\[
    \dim_{\mathbb{K}}\,\left(V_1 +
    V_2\right)
    =
    \dim_{\mathbb{K}}\,V_1+
    \dim_{\mathbb{K}}\,V_2
    \]
\end{theorem}
\begin{Proof}
设右式依次为$m$、$n$、$r$,
取$V_1\cap V_2$的一组基$\left\{\bm{\alpha}_1,\bm{\alpha}_2,\cdots,\bm{\alpha}_r
\right\}$,因为$V_1\cap V_2\subseteq V_1$,根据基扩张定理得到
$V_1$的一组基$\left\{\bm{\alpha}_1,\bm{\alpha}_2,\cdots,\bm{\alpha}_r
,\bm{\beta}_1,\bm{\beta}_2,\cdots,\bm{\beta}_{m-r}\right\}$,同理得到
$V_2$的一组基
$\left\{\bm{\alpha}_1,\bm{\alpha}_2,\cdots,\bm{\alpha}_r
,\bm{\gamma}_1,\bm{\gamma}_2,\cdots,\bm{\gamma}_{n-r}
\right\}$,所以只需要证明
\[
\left\{\bm{\alpha}_1,\bm{\alpha}_2,\cdots,\bm{\alpha}_r
,\bm{\beta}_1,\bm{\beta}_2,\cdots,\bm{\beta}_{m-r},
\bm{\gamma}_1,\bm{\gamma}_2,\cdots,\bm{\gamma}_{n-r}\right\}
\]
是$V_1+V_2$的基即可.

$\forall \bm{v}=\bm{v}_1+\bm{v}_2\in V_1+V_2 ,
\bm{v}_1\in V_1,\bm{v}_2\in V_2$,显然,
分成两部分分别考虑,容易得出$\bm{v}$是
$\left\{\bm{\alpha}_1,\bm{\alpha}_2,\cdots,\bm{\alpha}_r
,\bm{\beta}_1,\bm{\beta}_2,\cdots,\bm{\beta}_{m-r},
\bm{\gamma}_1,\bm{\gamma}_2,\cdots,\bm{\gamma}_{n-r}\right\}
$的线性组合.

再证线性无关.设
\begin{align*}
   &\lambda_1\bm{\alpha}_1+\lambda_2\bm{\alpha}_2+\cdots+\lambda_r
\bm{\alpha}_r
\\
+&\mu_1\bm{\beta}_1+\mu_2
\bm{\beta}_2+\cdots+\mu_{m-r}
\bm{\beta}_{m-r}\\
+&k_1
\bm{\gamma}_1+k_2
\bm{\gamma}_2+\cdots+k_{n-r}\bm{\gamma}_{n-r}
 =\bm{0}
\end{align*}
移项,有
\begin{align*}
   V_1\ni&\lambda_1\bm{\alpha}_1+\lambda_2\bm{\alpha}_2+\cdots+\lambda_r
\bm{\alpha}_r
+\mu_1\bm{\beta}_1+\mu_2
\bm{\beta}_2+\cdots+\mu_{m-r}
\bm{\beta}_{m-r}\\
=&-\left(k_1
\bm{\gamma}_1+k_2
\bm{\gamma}_2+\cdots+k_{n-r}\bm{\gamma}_{n-r}
\right)\in V_2
\end{align*}
故左右两边都属于$V_1\cap V_2.$不妨记作
\[
l_1\bm{\alpha}_1+l_2\bm{\alpha}_2+\cdots+l_r\bm{\alpha}_r
=-\left(k_1
\bm{\gamma}_1+k_2
\bm{\gamma}_2+\cdots+k_{n-r}\bm{\gamma}_{n-r}
\right)
\]
移项并由线性无关得$l_1=l_2=\cdots=l_r=k_1=k_2=\cdots=k_{n-r}=0.
$,代回又得
$\lambda_1=\lambda_2=\cdots=\lambda_r=\mu_1=\mu_2=\cdots=\mu_{m-r}=0.
$

证毕.
\end{Proof}
\end{formal}
\subsection{直和}
\begin{formal}
\begin{definition}[直和]\label{def:直和}
    设$V_1,V_2,\cdots,V_m$为$V$的子空间
    ,若$\forall 1
    \leqslant i\leqslant m$
    \[
    V_i \cap \left(
        V_1 + V_2+\cdots +\widehat{V_i}
        +\cdots +V_m
    \right)
     = 0
    \]则称和$V_1 
        +V_2+\cdots +V_m$
        为直接和,简称直和,记作\[
        V_1 \oplus V_2\oplus \cdots
        \oplus V_m
        \]
\end{definition}
\end{formal}
\begin{brown}
\begin{example}
    $m = 2$时,$V_1 + V_2$为直和等价于
    $V_1 \cap V_2 = 0$,比如$V = 
    \mathbb{R}^3 = V_1 \oplus 
    V_{23}$
\end{example}
\end{brown}
\begin{brown}
\begin{example}
    $m \geqslant 3$时$V_i$两两的交为$0$
    不能保证直和,比如
    \[
    V_1\oplus V_2 \oplus V_3
    \Longleftrightarrow
     \begin{cases*}
        V_1 \cap \left(V_2+V_3\right) = 0\\
        V_2 \cap \left(V_1+V_3\right) = 0 \\
        V_3 \cap \left(V_1+V_2\right) = 0
    \end{cases*}
    \]
\end{example}
\end{brown}
\begin{red}
\begin{remark}
    当子空间数大于2,子空间
    两两的交为零子空间并不能确定直和,
    直和的条件比两两的交为零子空间要严格得多.
\end{remark}
\end{red}
\begin{brown}
\begin{example}
    考虑笛卡尔坐标系,$V=\mathbb{R}^2$,
    $V_1:x$轴,$V_2:y$轴,$V_3:y=x$.即使两两的交为零子空间,
    但
    \[
    V_3\cap\left(V_1+V_2\right)=V_3\neq 0
    \]
\end{example}
\end{brown}
\begin{formal}
\begin{theorem}[直和的判定]\label{thm:直和的判定}
若$V_0 = V_1 +V_2+\cdots
+V_m$,则下列结论等价
\begin{enumerate}[label=\textup{(\arabic*)}]
    \item $V_0 =V_1\oplus V_2\oplus
    \cdots\oplus V_m$%为直和
    \item $\forall 2\leqslant i\leqslant m,
    V_i \cap \left(V_1 +\cdots
    +V_{i - 1}\right)
    = 0$
    \item $\dim\,\left( V_1 +V_2+\cdots
    +V_m\right)= \dim\, V_1 +\dim\,V_2+\cdots
    +\dim \,V_m$
    \item $V_1,V_2,\cdots,V_m$的各一组基可以
    拼成$V_0$的一组基
    \item $V_0$中向量用$V_1,V_2,\cdots,V_m
    $中向量和表出方式唯一,即若
    $\bm{\alpha}\in V_0 $且$
    \bm{\alpha} =
    \bm{v}_1+\bm{v}_2+\cdots
    +\bm{v}_m = 
    \bm{u}_1 + \bm{u}_2+\cdots
    + \bm{u}_m
    $其中$\bm{v}_i,\bm{u}_i
    \in V_i$则$\bm{v}_i = \bm{u}_i
    \left(
        1\leqslant i\leqslant m
    \right)$,称为分块表
    示唯一.事实上,这等价于
    零向量表示唯一.
\end{enumerate}
\end{theorem}
\begin{Proof}
\begin{enumerate}[label=\textup{(\Roman*)}]
\item $(1)\Longrightarrow (2)$
\[
V_i\cap \left(V_1+\cdots+V_{i-1}\right)\subseteq
V_i\cap \left(V_1+\cdots+V_{i-1}+V_{i+1}+\cdots+
V_m\right)=0
\]
故
\[
V_i \cap \left(V_1 +\cdots
    +V_{i - 1}\right)
    = 0\left(\forall 2\leqslant i\leqslant 
    m\right)
\]

\item $(2)\Longrightarrow (3)$对$V_m\cap\left(V_1+\cdots+V_{m-1}\right)=0$使用\cref{thm:子空间的维数公式}维数公式
得$\dim\,\left(V_1+\cdots+V_m\right)=\dim\,\left(V_1+\cdots+V_{m-1}\right)+\dim\,V_m$,一直如此便得到
\[
\dim\,\left( V_1 +V_2+\cdots
    +V_m\right)= \dim\, V_1 +\dim\,V_2+\cdots
    +\dim \,V_m
\]

\item $(3)\Longrightarrow(4)$设
$\dim\,V_i=n_i$,维数$\dim\,V_0=n_1+\cdots+n_m
=n$.取$V_i$的一组基
$\left\{\bm{e}_{i1},\bm{e}_{i2},\cdots,
\bm{e}_{in_i}\right\}\left(1\leqslant i \leqslant 
m\right)$,拼在一起即
\[
\left\{
    \bm{e}_{11},\bm{e}_{12},\cdots,\bm{e}_{1n_1},
    \bm{e}_{21},\bm{e}_{22},\cdots,\bm{e}_{2n_2},
    \cdots,
    \bm{e}_{m1},\bm{e}_{m2},\cdots,\bm{e}_{mn_m}
\right\}
\]
共$n_1+n_2+\cdots+n_m=n$个.又$\dim\,V_0=n.$

$\forall \bm{\alpha}=\bm{\alpha}_1+\cdots+\bm{\alpha}_m\in V_0=
V_1+\cdots+V_m$,其中$\bm{\alpha}_i\in V_i.$
设$\bm{\alpha}_i=
\lambda_{i1}\bm{e}_{i1}+\lambda_{i2}\bm{e}_{i2}+\cdots+\lambda_{in_i}
\bm{e}_{in_i}\left(1\leqslant i \leqslant m\right)$.容易证明线性组合,
以此判定基.

\item $(4)\Longrightarrow (5)$,我们先证其与推论等价,
 由$(5)$推出推论显然；反之,
 \[
 \bm{0}=\left(\bm{v}_1-\bm{u}_1\right)+\left(
    \bm{v}_2-\bm{u}_2
 \right)+\cdots+\left(\bm{v}_m-\bm{u}_m\right)\in V_1+V_2+\cdots+V_m
 \]
 由零向量表示唯一立得$\bm{v}_i=\bm{u}_i.$
 于是只需要证明$(4)$推得$(5)$推论.

设$\bm{0}=\bm{v}_1+\bm{v}_2+\cdots+\bm{v}_m$,
其中$\bm{v}_i\in V_i.$因为
$\bm{v}_i=
\lambda_{i1}\bm{e}_{i1}+\lambda_{i2}\bm{e}_{i2}+\cdots+\lambda_{in_i}
\bm{e}_{in_i}\left(
    1\leqslant i\leqslant m
\right)$.故
$\displaystyle
\bm{0}=\sum_{i=1}^{m}\left(
    \lambda_{i1}\bm{e}_{i1}+\lambda_{i2}\bm{e}_{i2}+\cdots+\lambda_{in_i}
\bm{e}_{in_i}
\right)$,由基的线性无关性得
\[
\lambda_{ij}=0,\forall 1\leqslant i\leqslant m,1\leqslant j\leqslant n_i
\]
即$\forall 1\leqslant i\leqslant m,\bm{v}_i=\bm{0}$.

\item 最后为该推论推得$(1)$.任取$
\bm{\alpha}\in V_i\cap \left(V_1+\cdots+V_{i-1}+V_{i+1}+\cdots+
V_m\right)$,
即$\bm{\alpha}=\bm{\alpha}_1+\cdots+\bm{\alpha}_{i-1}+\bm{\alpha}_{i+1}+\cdots+\bm{\alpha}_m
$.即有对零向量分解
\[
\bm{0}=\bm{\alpha}_1+\cdots
+\bm{\alpha}_{i-1}+\left(-\bm{\alpha}
\right)+\bm{\alpha}_{i+1}+\cdots+\bm{\alpha}_m
\]
则$\bm{\alpha}_1=\cdots=\bm{\alpha}_{i-1}=\bm{\alpha}=
\bm{\alpha}_{i+1}=\cdots=\bm{\alpha}_m=\bm{0}.$
即$\bm{\alpha}=0$,于是
\[
V_i\cap \left(V_1+\cdots+V_{i-1}+V_{i+1}+\cdots+
V_m\right)=0\left(1\leqslant i\leqslant 
m\right)
\qedhere
\]
\end{enumerate}
\end{Proof}
\end{formal}
\newpage
\section{线性方程组的解}
\subsection{判定}
\begin{formal}
\begin{theorem}[判定定理]\label{thm:判定定理}
    线性方程组有解与否及是否有非平凡解的判定.
    \[
    \bm{A}\bm{x} = \bm{\beta}
    \Longleftrightarrow 
    \begin{cases*}
        a_{11}x_1+a_{12}x_2+\cdots+a_{1n}x_n = b_1 \\
        a_{21}x_1+a_{22}x_2+\cdots+a_{2n}x_n=b_2\\
        \qquad\cdots\cdots\cdots\cdots\cdots
        \cdots\cdots \\
        a_{m1}x_1+a_{m2}x_2+\cdots+a_{mn}x_n=b_m
    \end{cases*}
    \]
    其增广矩阵
    \[
    \widetilde{\bm{A}}
    =\left (
    \begin{array}{c:c}
        \begin{matrix}
            \bm{A}
        \end{matrix}
        &
        \begin{matrix}
            \bm{\beta}
        \end{matrix}
    \end{array}
    \right )
    =\left (
        \begin{array}{c:c}
            \begin{matrix}
                a_{11} & a_{12} & \cdots & a_{1n} \\
                a_{21} & a_{22} & \cdots & a_{2n} \\
                \vdots & \vdots &  & \vdots \\
                a_{m1} & a_{m2} & \cdots & a_{mn}
            \end{matrix}
            &
            \begin{matrix}
                b_1 \\
                b_2 \\
                \vdots
                \\
                b_m
            \end{matrix}
        \end{array}
    \right )
    \]

    \begin{enumerate}[label=\textup{(\arabic*)}]
      \item $ \mathrm{r}\left(\bm{A}\right)  
      =\mathrm{r}\left(\widetilde{\bm{A}}
      \right) = n\Longrightarrow 
        $有唯一一组解
      \item $\mathrm{r}\left(\bm{A}\right)  
      =\mathrm{r}\left(\widetilde{\bm{A}}
      \right) < n \Longrightarrow $
        有无穷多组解
      \item $\mathrm{r}\left(\bm{A}\right)
      \neq \mathrm{r}\left(
        \widetilde{\bm{A}}
      \right)($
        事实上,可以断言
      $\mathrm{r}\left(
        \widetilde{\bm{A}
      }\right) = 
      \mathrm{r}\left(\bm{A}\right) + 1
      $)$\Longrightarrow $
      无解
    \end{enumerate}
\end{theorem}
\begin{Proof}
先证有解当且仅当
$ \mathrm{r}\left(\bm{A}\right)  
      =\mathrm{r}\left(\widetilde{\bm{A}}
      \right)$.作列分块$\bm{A}=\left(\bm{\alpha}_1,
      \bm{\alpha}_2,\cdots,\bm{\alpha}_n
      \right)$,于是方程化为
      \[
      x_1\bm{\alpha}_1+x_2\bm{\alpha}_2+\cdots+x_n\bm{\alpha}_n
      =\bm{\beta}
      \]
因此方程有解等价于$\bm{\beta}$是$\bm{\alpha}_1,
\bm{\alpha}_2,\cdots,\bm{\alpha}_n$的线
性组合.

先证必要性,设方程有解,则
$\bm{\beta}$是$\bm{\alpha}_1,
\bm{\alpha}_2,\cdots,\bm{\alpha}_n$的线
性组合.设$\bm{A}$的列向量的极大无关组为
$\left\{\bm{\alpha}_{i_1},\bm{\alpha}_{i_2}
,\cdots,\bm{\alpha}_{i_r}\right\}$,
即$\mathrm{r}\left(\bm{A}\right)=r$.显然
$\bm{\beta}$是$\bm{\alpha}_{i_1},
\bm{\alpha}_{i_2},\cdots,\bm{\alpha}_{i_r}
$的线
性组合.于是
$\left\{\bm{\alpha}_{i_1},
\bm{\alpha}_{i_2},\cdots,\bm{\alpha}_{i_r}
\right\}$也是$\widetilde{\bm{A}}$的
列向量的极大无关组.于是$\mathrm{r}\left(\widetilde{\bm{A}}\right)=r=\mathrm{r}\left(\bm{A}\right).
$

再证充分性,设$\mathrm{r}\left(\widetilde{\bm{A}}\right)=
\mathrm{r}\left(\bm{A}\right)=r$,并且
$\left\{\bm{\alpha}_{i_1},\bm{\alpha}_{i_2},
\cdots,\bm{\alpha}_{i_r}\right\}$是$\bm{A}$的列向量的极大无关组,
即也是$\widetilde{\bm{A}}$的$
r$个线性无关的向量,
即是$\widetilde{\bm{A}}$的列向量
的极大无关组.于是
$\bm{\beta}$是$\bm{\alpha}_{i_1},
\bm{\alpha}_{i_2},\cdots,\bm{\alpha}_{i_r}
$的线
性组合.

有解条件证毕.

若$\mathrm{r}\left(\bm{A}\right)=\mathrm{r}\left(
    \widetilde{\bm{A}}
\right)=n$,则$\left(\bm{\alpha}_1,
\bm{\alpha}_2,\cdots,\bm{\alpha}_n\right)$线
性无关.又因为线性表示唯一,
因此$\left(x_1,x_2,\cdots,x_n\right)$唯一确定.

若$\mathrm{r}\left(\bm{A}\right)=\mathrm{r}\left(
    \widetilde{\bm{A}}
\right)<n$,则$\left(\bm{\alpha}_1,
\bm{\alpha}_2,\cdots,\bm{\alpha}_n\right)$线
性相关即存在不全为零的$c_1,c_2,\cdots,c_n\in \mathbb{K}$使得
\begin{align*}
    \bm{0}=c_1\bm{\alpha}_1
+c_2\bm{\alpha}_2+\cdots+
c_n\bm{\alpha}_n
\end{align*}
又因为方程有解等价于存在$k_1,k_2,\cdots,k_n\in \mathbb{K}$使得
\[
\bm{\beta}=k_1\bm{\alpha}_1+k_2\bm{\alpha}_2+\cdots+k_n\bm{\alpha}_n
\]
两式组合得
\[
\bm{\beta}=
\left(k_1+kc_1\right)\bm{\alpha}_1+\left(
    k_2+kc_2
\right)\bm{\alpha}_2+\cdots+\left(k_n+kc_n\right)
\bm{\alpha}_n
\]
即$\displaystyle
x_i=k_i+kc_i,\forall k\in \mathbb{K}$,即有无
穷多组解.

证毕.
\end{Proof}
\end{formal}
\subsection{解的结构与表达}
当线性方程组具有无穷多组解时,
一个重要的问题
就是如何用人有限多组解来表达或者说生成
这无穷多组解.
\begin{green}
\begin{lemma}[相伴方程组的解]\label{lem:相伴方程组的解}
    设$\bm{\gamma}$是$\bm{Ax}=\bm{\beta}$
    的一个特解,则$\bm{\alpha}$是$
    \bm{Ax}=\bm{\beta}$的解等价于
    $\bm{\alpha}-\bm{\gamma}$为相伴的
    齐次线性方程组$\bm{Ax}=\bm{0}$的解.
\end{lemma}
\begin{Proof}
证明略.
\end{Proof}
\end{green}
\begin{green}
\begin{lemma}[解空间]\label{lem:解空间}
    考虑相伴的齐次线性方程组$\bm{Ax}=\bm{0}$,其至少有零解.
    令
    \[
    V_{\bm{A}}=\left\{\bm{x}\in \mathbb{K}^n\big|\bm{Ax}=\bm{0}\right\}
    \]
    即方程的解集,那么断言这是$\mathbb{K}^n$的子空间,称为
    该齐次线性方程组的解空间.
\end{lemma}
\end{green}
\begin{formal}
\begin{theorem}[基础解系]\label{thm:基础解系}
    设$\mathrm{r}\left(\bm{A}
    \right)=r$,则解空间
    $V_{\bm{A}} = 
    \left\{
        \bm{x} \in \mathbb{K}^n\mid
        \bm{Ax} =\bm{0}
    \right\}$是
    $\mathbb{K}^n$的一个$n  -r$维
    子空间即
    \[
    \dim_{\mathbb{K}}\,V_{\bm{A}} + 
    \mathrm{r}\left(\bm{A}\right)
    = n
    \]
    从而存在一组基
    $\left\{\bm{\eta}_1,\bm{\eta}_2,\cdots
    ,\bm{\eta}_{n - r}\right\}$
    使得$\bm{Ax}=\bm{0}$的解都是
    $\bm{\eta}_1,\bm{\eta}_2,\cdots,
    \bm{\eta}_{n - r}$的线性组合,
    这组基$\left\{\bm{\eta}_1,\bm{\eta}_2,\cdots
    ,\bm{\eta}_{n - r}\right\}$
    （齐次线性方程组解空间的基）称为
    基础解系.
\end{theorem}
\begin{Proof}
首先考虑高斯消去法与初等行变换,我们知道,
对线性方程组进行同解基础上的化简等价于对
系数矩阵$\bm{A}$做初等行变换.
此外,列对换仅仅改变了未定元的
位置,因此是允许的.

由行变换,可将$\bm{A}$的行向量的极大无关组换至
前$r$行.不妨设
\[
\bm{A}=\begin{pmatrix}
    \bm{\alpha}_1\\
    \bm{\alpha}_2\\
    \vdots\\
    \bm{\alpha}_r\\
    \bm{\alpha}_{r+1}\\
    \vdots\\
    \bm{\alpha}_m
\end{pmatrix}
\]
其中$\left\{\bm{\alpha}_1,
\bm{\alpha}_2,
\cdots,\bm{\alpha}_r\right\}$为
行向量的极大无关组.于是有
\[
\bm{A}\longrightarrow
\begin{pmatrix}
    \bm{\alpha}_1\\
    \bm{\alpha}_2\\
    \vdots\\
    \bm{\alpha}_r\\
    \bm{0}\\
    \vdots\\
    \bm{0}
\end{pmatrix}
\]
令\[
\bm{A}_1=\begin{pmatrix}
    \bm{\alpha}_1\\\bm{\alpha}_2\\\vdots\\\bm{\alpha}_r
\end{pmatrix}\Longrightarrow
\mathrm{r}\left(\bm{A}_1\right)=r
\]
在允许列对换（交换变量位置）的情况下,不妨设$\bm{A}_1$的列向量的极大无关组
为前$r$列,则
$\bm{A}_1=\left(\bm{B}_1^{r\times r}
,\bm{B}_2^{r\times \left(n-r\right)}\right)$,
则$\mathrm{r}\left(\bm{B}_1\right)=r.$
故$\bm{B}_1$非异,其可通过行变换成为单位阵,此时
$\bm{A}\longrightarrow \left(\bm{I}_r,
\bm{C}\right)$.

总之,在初等行变换和列对换后$\bm{A}$化为
\[
\begin{bmatrix}
    \bm{I}_r&\bm{C}\\\bm{O}&\bm{O}
\end{bmatrix}
,\bm{C}=\left(c_{ij}\right)_{r\times \left(n-r\right)}
\]
从而所给齐次线性方程组与下面的方程同解
\[
\begin{cases*}
x_1+\qquad \qquad +c_{1,r+1}x_{r+1}+\cdots+c_{1,n}x_n=0\\
\qquad x_2+\qquad +c_{2,r+1}x_{r+1}+\cdots+c_{2,n}x_n=0\\
\qquad\qquad\cdots\cdots\cdots\cdots\cdots\cdots\cdots\cdots\\
\qquad\qquad \quad x_r+c_{r,r+1}x_{r+1}+\cdots+c_{r,n}x_n=0
\end{cases*}
\]
取$x_{r+1}=1$,$x_{r+2}=\cdots=x_n=0$得一个解
\[
\bm{\eta}_1=\begin{pmatrix}
    -c_{1,r+1}\\
    -c_{2,r+1}\\
    \vdots\\
    -c_{r,r+1}\\
    1\\
    0\\
    \vdots\\
    0
\end{pmatrix}
\]
随后令
$x_{r+2}=1$,其他为零,如此下去$\cdots$,得到
$\bm{\eta}_1,\bm{\eta}_2,\cdots,\bm{\eta}_{n-r}
$
,并断言这就是该方程解空间的一组基.

一方面,它们显然线性无关；另一方面,任取上述方程一解
$\bm{\eta}=\left(a_1,a_2,\cdots,a_n\right)
'$
,那么
\[
\bm{\eta}=\begin{pmatrix}
    a_1\\a_2\\\vdots\\a_n
\end{pmatrix}
=\begin{pmatrix}
    -c_{1,r+1}a_{r+1}-\cdots-c_{1,n}a_n\\
    -c_{2,r+1}a_{r+1}-\cdots-c_{2,n}a_n\\
    \vdots\\
    -c_{r,r+1}a_{r+1}-\cdots-c_{r,n}a_n\\
    a_{r+1}\\
    \vdots\\
    a_n
\end{pmatrix}=a_{r+1}\bm{\eta}_1+\cdots+a_n\bm{\eta}
_{n-r}
\]
线性组合亦成立.故这是的该方程解空间的一
组基.
同时也是齐次线性方程组的解空间$V_{\bm{A}}$的一组基.

于是
\[
\dim\,V_{\bm{A}}=n-r=n-\mathrm{r}\left(\bm{A}
\right)
\qedhere
\]
\end{Proof}
\end{formal}
\begin{formal}
\begin{theorem}[结构定理]\label{thm:结构定理}
    设$\mathrm{r}\left(\bm{A}\right)
    =\mathrm{r}\left(\widetilde{\bm{A}
    }\right) = r$,而$\bm{\gamma}$是$\bm{Ax}=\bm{\beta}$
    的一个特解,$\left\{\bm{\eta}_1,\bm{\eta}_2,\cdots
    ,\bm{\eta}_{n - r}\right\}$是相伴
    齐次线性方程组$\bm{Ax}=\bm{0}$
    的一个基础解系,则$
    \bm{Ax}=\bm{\beta}$的通解为
    \[
        \bm{\gamma}+
        k_1\bm{\eta}_1+k_2\bm{\eta}_2+\cdots
        +k_{n - r}\bm{\eta}_{n - r}
    \]
    其中$k_i \in \mathbb{K}$.
\end{theorem}
\begin{Proof}
考虑$\bm{\alpha}-\bm{\gamma}$用基础解
系线性组合立得.
\end{Proof}
\end{formal}
\begin{red}
\begin{remark}
    强调数域$\mathbb{K}$求解.
\end{remark}
\end{red}
\begin{red}
\begin{remark}
    求解过程中,仅允许初等行变换和第三类初等列变换即列对换,
    其它变换不允许使用.

    注意,是列对换.
\end{remark}
\end{red}
\subsection{求解$\bm{Ax}=\bm{\beta}$}
\begin{formal}
\begin{theorem}[求解$\bm{Ax}=\bm{\beta}$的方法]\label{thm:求解Ax=b的方法}
一般地,求解方程
\[
\begin{cases*}
    a_{11}x_1+a_{12}x_2+\cdots+a_{1n}x_n = b_1 \\
    a_{21}x_1+a_{22}x_2+\cdots+a_{2n}x_n=b_2\\
    \qquad\cdots\cdots\cdots\cdots\cdots
    \cdots\cdots \\
    a_{m1}x_1+a_{m2}x_2+\cdots+a_{mn}x_n=b_m
\end{cases*}\Longleftrightarrow
\bm{Ax}=\bm{\beta}
\]
的方法如下：

$1.$施加行变换将$\widetilde{\bm{A}}
    =\left (
    \begin{array}{c:c}
        \begin{matrix}
            \bm{A}
        \end{matrix}
        &
        \begin{matrix}
            \bm{\beta}
        \end{matrix}
    \end{array}
    \right )$化为阶梯形,判断$\mathrm{r}\left(\bm{A}
    \right)$与$\mathrm{r}\left(\widetilde{
        \bm{A}
    }\right)$的关系来确定解的存在性.

$2.$继续对$\widetilde{\bm{A}}$作
初等行变换和
列对换（不允许采用其他变换
,不允许在列对换时
操作常数列,列对换仅仅变换了
未定元的次序,后面需要换回来）
化为如下解方程组的标准型形式
\[
\bm{D} =
\left [
    \begin{array}{c:c}
        \begin{matrix}
            \bm{I}_r & \bm{C}_{r
            \times \left(n - r\right)}\\
            \bm{O} & \bm{O}
        \end{matrix}
        &
        \begin{matrix}
            \bm{\gamma}\\
            \bm{O}
        \end{matrix}
    \end{array}
\right ]
\]
得到特解
\[
\begin{bmatrix}
    \bm{\gamma}\\
    \bm{O}
\end{bmatrix}
\]和基础解系
$\bm{\eta}_1,\bm{\eta}_2,\cdots,
\bm{\eta}_{n - r}$,其中
\[
\bm{\eta}_i
=\left ( 
    \begin{array}{c}
        -c_{1,r+i}\\
        -c_{2,r+i}\\
        \vdots \\
        -c_{r,r+i}\\
        \hdashline
        0\\
        \vdots \\
        1\\
        \vdots\\
        0
    \end{array}
\right )
\]
其中,$c_{i,r+i}$是$\bm{C}_{
    r\times\left(n - r\right)
}$的第$\left(i,i\right)$元或者说
$\bm{D}$的$\left(i,r+i\right)$元；
元$1$位于$-c_{r,r+i}$下第$i$个位
置或者说是总的第$r + i$个位置.

$3.$根据列对换的情况调整特解
$\bm{\gamma}$和基础解系
$\left\{\bm{\eta}_1,\bm{\eta}_2,\cdots,\bm{\eta}_{n - r}
\right\}$
的分量得到正确的特解
$\bm{\delta }$和基础解系
$\left\{\bm{\xi}_1,\bm{\xi}_2,\cdots
,\bm{\xi}_{n - r}\right\}$,最后组合为
通解
\[
\bm{\delta}+k_1\bm{\xi}_1+k_2\bm{\xi}_2
+\cdots+k_{n - r}\bm{\xi}_{n - r}
\]
其中$k_i \in \mathbb{K},\forall 1\leqslant i\leqslant n-r.$
\end{theorem}
\end{formal}
\begin{green}
\begin{corollary}[对解的一些刻画]\label{cor:对解的一些刻画}
    设$\bm{Ax} = \bm{\beta}\left(
        \bm{\beta} \neq \bm{0}
    \right)$的特解为$\bm{\gamma}$,
    $\bm{Ax} = \bm{0}$的基础解系为
    $\left\{\bm{\eta}_1,\bm{\eta}_2,\cdots
    ,\bm{\eta}_{n - r}\right\}$,则
    \begin{enumerate}[label=\textup{(\arabic*)}]
        \item $\bm{\gamma},\bm{\gamma}+
        \bm{\eta}_1,\bm{\gamma}+\bm{\eta}_2,\cdots,
        \bm{\gamma}+\bm{\eta}_{n - r}
        $线性无关
        \item $\bm{Ax}=\bm{\beta}$的解
        必符合如下形式：
        $c_0\bm{\gamma}+c_1\left(
            \bm{\gamma}+\bm{\eta}_1
        \right)+c_2\left(
            \bm{\gamma}+\bm{\eta}_2
        \right)+\cdots+
        c_{n - r}\left(\bm{\gamma}+
        \bm{\eta}_{n - r}\right)
        $并且有
        $c_0+c_1+c_2+\cdots+c_{n - r}=1$.
        反之,这样形式的
        向量一定是方程的解.
    \end{enumerate}
\end{corollary}
\begin{Proof}
$(1)$设
\[
\lambda_0\bm{\gamma}
+\lambda_1\left(
            \bm{\gamma}+\bm{\eta}_1
        \right)+\cdots+
        \lambda_{n - r}\left(\bm{\gamma}+
        \bm{\eta}_{n - r}\right)=\bm{0}
\]
即
\[
\left(\sum_{i=0}^{n-r}
\lambda_i\right)\bm{\gamma}+\lambda_1\bm{\eta}_1+\cdots
+\lambda_{n-r}\bm{\eta}_{n-r}=
\bm{0}
\]
用$\bm{A}$左乘,即得
\[
\left(\sum_{i=0}^{n-r}\lambda_{i}\right)\bm{\beta}=\bm{0}
\]
故$\displaystyle
\sum_{i=0}^{n-r}\lambda_i=0$,反代即得
\[
\lambda_1\bm{\eta}_1+\cdots
+\lambda_{n-r}\bm{\eta}_{n-r}=
\bm{0}
\]
又因为线性无关,
故$\lambda_1=\cdots=\lambda_{n-r}=0$,易推知$\lambda_0=0$.
证毕.

$(2)$由结构定理,任一解$\bm{\alpha}=
\bm{\gamma}+
        k_1\bm{\eta}_1+\cdots
        +k_{n - r}\bm{\eta}_{n - r},k_i\in \mathbb{K}.
$
显然有
\[
\bm{\alpha}=\left(1-k_1-k_2-\cdots-k_{n-r}\right)
\bm{\gamma}+k_1\left(\bm{\gamma}+\bm{\eta}_1\right)
+k_2\left(\bm{\gamma}+\bm{\eta}_2\right)+
\cdots+k_{n-r}\left(\bm{\gamma}+\bm{\eta}_{n-r}\right)
\qedhere
\]
\end{Proof}
\end{green}
\subsection{线性方程组求解的应用}
\begin{formal}
\begin{proposition}[线性方程组求解的相关应用]\label{prop:线性方程组求解的相关应用}
    核心是
\[
\dim_{\mathbb{K}}\,V_{\bm{A}} +
\mathrm{r}\left(\bm{A}\right)=n
\]
其中$n$为未定元个数.
主要可用于：
\begin{enumerate}[label=\textup{\arabic*}]
    \item $n$阶方阵$\bm{A}$可逆等价于
    $\bm{Ax}=\bm{0}$仅有零解.
    \item 利用$\mathrm{r}\left(\bm{A}
    \right)$求$V_{\bm{A}}$.
    \item 利用$V_{\bm{A}}$求$\mathrm{r}\left(
        \bm{A}\right)$.
\end{enumerate}
\end{proposition}
\end{formal}
\begin{brown}
\begin{example}
    已知$\bm{A}^2-\bm{A}-3\bm{I}_n=\bm{O}$,求证：
    $\bm{A}-2\bm{I}_n$非异.
\end{example}
\begin{proof}
    只要证方程
    \[
    \left(\bm{A}-2\bm{I}_n\right)\bm{x}=\bm{0}
    \]
    仅有零解.设一解$\bm{x}_0$,则
    $\bm{Ax}_0=2\bm{x}_0$,于是
    $\bm{A}^2\bm{x}_0=2\bm{Ax}_0=4\bm{x}_0$,
    故
    \[
    \left(\bm{A}^2-\bm{A}-3\bm{I}_n\right)\bm{x}_0=-\bm{x}_0=\bm{0}\Longrightarrow \bm{x}_0
    =\bm{0}
    \qedhere
    \]
\end{proof}
\end{brown}
\begin{brown}
\begin{example}
    设$\lambda_1,\lambda_2,\cdots,\lambda_n$是
    数域$\mathbb{K}$中不同的数且有$1\leqslant k \leqslant n-1$.
    设
    \begin{flalign}
        \tag{1}
        &\begin{cases*}
        x_1+x_2+\cdots+x_n=0\\
        \lambda_1x_1+\lambda_2x_2+\cdots+\lambda_nx_n=0\\
        \qquad\cdots\cdots\cdots\cdots\cdots\\
        \lambda_{1}^{k-1}x_1+\lambda_2^{k-1}x_2+\cdots+\lambda_n^{k-1}x_n=0
        \end{cases*}\\
        \tag{2}&
        \begin{cases*}
        x_1^{k}+x_2^k+\cdots+x^k_n=0\\
        \lambda_1^{k+1}x_1+\lambda_2^{k+1}x_2+\cdots+\lambda_n^{k+1}x_n=0\\
        \qquad\cdots\cdots\cdots\cdots\cdots\\
        \lambda_{1}^{n-1}x_1+\lambda_2^{n-1}x_2+\cdots+\lambda_n^{n-1}x_n=0
        \end{cases*}&
    \end{flalign}
设解空间分别为$V_1$、$V_2$,求证：$\mathbb{K}^n=V_1\oplus V_2.$
\end{example}
\begin{Proof}
考虑$V_1\cap V_2$,即$(1)$、$(2)$联立得到方程组的解空间
\[
\tag{3}
        \begin{cases*}
        x_1+x_2+\cdots+x_n=0\\
        \lambda_1x_1+\lambda_2x_2+\cdots+\lambda_nx_n=0\\
        \qquad\cdots\cdots\cdots\cdots\cdots\\
        \lambda_{1}^{k-1}x_1+\lambda_2^{k-1}x_2+\cdots+\lambda_n^{k-1}x_n=0\\
     x_1^{k}+x_2^k+\cdots+x^k_n=0\\
        \lambda_1^{k+1}x_1+\lambda_2^{k+1}x_2+\cdots+\lambda_n^{k+1}x_n=0\\
        \qquad\cdots\cdots\cdots\cdots\cdots\\
        \lambda_{1}^{n-1}x_1+\lambda_2^{n-1}x_2+\cdots+\lambda_n^{n-1}x_n=0    
    \end{cases*}
\]
其系数矩阵的行列式为Vandermonde行列式,值为
\[
\prod_{1\leqslant i<j\leqslant n}
\left(\lambda_j-\lambda_i\right)
\neq 0
\]
故系数矩阵满秩.于是$V_1\cap V_2=0.$
不妨将三个系数矩阵记作
\newpage
\[
\bm{A}=\begin{pmatrix}
    \bm{A}_1\\\bm{A}_2
\end{pmatrix}
\]
且$\mathrm{r}\left(\bm{A}\right)=n$,则
$\mathrm{r}\left(\bm{A}_1\right)=k,\mathrm{r}\left(\bm{A}_2\right)=n-k.
$于是$\dim\,V_1=n-\mathrm{r}\left(\bm{A}_1\right)=n-k,
\dim\,V_2=n-\mathrm{r}\left(\bm{A}_2\right)=k
.$故$\dim\,\left(V_1\oplus V_2\right)=n=\dim\,\mathbb{K}^n.
$
\end{Proof}
\end{brown}
\begin{brown}
\begin{example}
    设$\bm{A}\in M_{m\times n}\left(\mathbb{R}\right)$,求证：
    $\mathrm{rank}\left(\bm{AA}'\right)=\mathrm{rank}\left(\bm{A}'\bm{A}\right)
    =\mathrm{rank}\left(\bm{A}\right).$
\end{example}
\begin{Proof}
设方程$\bm{Ax}=\bm{0}$的解空间为$V_{\bm{A}}$,方程
$\bm{A}'\bm{Ax}=\bm{0}$的解空间为$V_{\bm{A}'\bm{A}}
$.显然
\newpage
\[
V_{\bm{A}}\subseteq
 V_{\bm{A}'\bm{A}}
 \]
 任取$\bm{x}_0\in V_{\bm{A}'\bm{A}}$,此时
 $\bm{x}_0\in \mathbb{R}^n$且$\bm{A}'\bm{Ax}_0=\bm{0}.
 $令$\bm{Ax}_0=\left(a_1,a_2,\cdots,a_m\right)'\in \mathbb{R}^m$.
 对$\bm{A}'\bm{Ax}_0=\bm{0}$左乘$\bm{x}_0'$得
 \[
\left(\bm{Ax}_0\right)'\left(\bm{Ax}_0\right)
=\bm{0}
 \]
 即$\displaystyle
 \sum_{i=1}^{m}a_i^2=0.$故$\forall 1\leqslant i\leqslant m,a_i=0$即
 $\bm{Ax}_0=\bm{0}.$则
 \[
V_{\bm{A}}\supseteq
 V_{\bm{A}'\bm{A}}
 \]
 故$
V_{\bm{A}}=
 V_{\bm{A}'\bm{A}}
 $,利用维数公式即得.
\end{Proof}
\end{brown}